% ============================================================
%  AI INFERENCE ENGINEERING
%  Building Fast, Scalable, and Production-Ready AI Systems
%  Full Book — LaTeX Source
%  Updated: May 2026 — covering H100/Blackwell, vLLM, TRT-LLM
% ============================================================

\documentclass[11pt, twoside, openright]{book}

% ─── Page geometry ───────────────────────────────────────────
\usepackage[
  paperwidth=7.5in, paperheight=9.25in,
  top=1in, bottom=1in,
  inner=1.1in, outer=0.85in,
  includeheadfoot
]{geometry}

% ─── Core packages ───────────────────────────────────────────
\usepackage[T1]{fontenc}
\usepackage[utf8]{inputenc}
\usepackage{lmodern}
\usepackage{microtype}
\usepackage{amsmath, amssymb}
\usepackage{graphicx}
\usepackage{xcolor}
\usepackage{booktabs}
\usepackage{longtable}
\usepackage{array}
\usepackage{enumitem}
\usepackage{fancyhdr}
\usepackage{titlesec}
\usepackage{imakeidx}        % index
\usepackage{mdframed}
\usepackage{tcolorbox}
\tcbuselibrary{skins, breakable, listings}
\usepackage{tikz}
\usetikzlibrary{shapes.geometric, arrows.meta, positioning, fit,
                decorations.pathreplacing, calc, backgrounds}
\usepackage{listings}
\usepackage{hyperref}
\usepackage{caption}
\usepackage{subcaption}
\usepackage{setspace}
\usepackage{emptypage}

% ─── Hyperref setup ──────────────────────────────────────────
\hypersetup{
  colorlinks  = true,
  linkcolor   = NavyBlue,
  citecolor   = NavyBlue,
  urlcolor    = DarkBlue,
  pdftitle    = {AI Inference Engineering},
  pdfauthor   = {AI Inference Engineering Book},
  pdfsubject  = {AI Systems, GPU Programming, Production ML},
  bookmarksnumbered = true,
  bookmarksopen     = true
}

% ─── Colors ──────────────────────────────────────────────────
\definecolor{NavyBlue}{RGB}{0, 51, 102}
\definecolor{DarkBlue}{RGB}{0, 70, 140}
\definecolor{AccentOrange}{RGB}{220, 90, 0}
\definecolor{AccentGreen}{RGB}{0, 128, 64}
\definecolor{LightGray}{RGB}{245, 245, 245}
\definecolor{MidGray}{RGB}{130, 130, 130}
\definecolor{DarkGray}{RGB}{50, 50, 50}
\definecolor{CodeBg}{RGB}{248, 248, 252}
\definecolor{CodeBorder}{RGB}{180, 190, 220}
\definecolor{TipBg}{RGB}{230, 245, 230}
\definecolor{TipBorder}{RGB}{0, 150, 60}
\definecolor{WarnBg}{RGB}{255, 245, 220}
\definecolor{WarnBorder}{RGB}{200, 130, 0}
\definecolor{NoteBg}{RGB}{225, 235, 255}
\definecolor{NoteBorder}{RGB}{60, 100, 200}
\definecolor{KeywordColor}{RGB}{0, 0, 180}
\definecolor{CommentColor}{RGB}{100, 120, 100}
\definecolor{StringColor}{RGB}{160, 40, 0}
\definecolor{NumberColor}{RGB}{140, 0, 140}

% ─── Index ───────────────────────────────────────────────────
\makeindex[columns=2, title=Index, intoc]

% ─── Chapter title style ─────────────────────────────────────
\titleformat{\chapter}[display]
  {\normalfont\bfseries\color{NavyBlue}}
  {\small\color{MidGray}\MakeUppercase{\chaptertitlename}\ \thechapter}
  {0.5em}
  {\huge}
  [\vspace{0.3em}{\color{AccentOrange}\titlerule[2pt]}]

\titlespacing*{\chapter}{0pt}{-10pt}{30pt}

\titleformat{\section}
  {\normalfont\Large\bfseries\color{NavyBlue}}
  {\thesection}{1em}{}

\titleformat{\subsection}
  {\normalfont\large\bfseries\color{DarkGray}}
  {\thesubsection}{1em}{}

\titleformat{\subsubsection}
  {\normalfont\normalsize\bfseries\color{DarkGray}}
  {\thesubsubsection}{1em}{}

% ─── Headers/Footers ─────────────────────────────────────────
\pagestyle{fancy}
\fancyhf{}
\fancyhead[LE]{\small\color{MidGray}\textit{\leftmark}}
\fancyhead[RO]{\small\color{MidGray}\textit{\rightmark}}
\fancyfoot[LE,RO]{\small\color{MidGray}\thepage}
\fancyfoot[LO,RE]{\small\color{MidGray}\textit{AI Inference Engineering}}
\renewcommand{\headrulewidth}{0.4pt}
\renewcommand{\footrulewidth}{0.4pt}

% ─── Listings (code) style ───────────────────────────────────
\lstdefinestyle{pythonstyle}{
  language=Python,
  backgroundcolor=\color{CodeBg},
  basicstyle=\ttfamily\small,
  keywordstyle=\color{KeywordColor}\bfseries,
  commentstyle=\color{CommentColor}\itshape,
  stringstyle=\color{StringColor},
  numberstyle=\tiny\color{MidGray},
  numbers=left,
  numbersep=8pt,
  stepnumber=1,
  breaklines=true,
  breakatwhitespace=true,
  showstringspaces=false,
  frame=single,
  framesep=4pt,
  rulecolor=\color{CodeBorder},
  tabsize=4,
  captionpos=b,
  morekeywords={import,from,as,def,class,return,yield,async,await,
                with,pass,raise,True,False,None},
}

\lstdefinestyle{bashstyle}{
  language=bash,
  backgroundcolor=\color{LightGray},
  basicstyle=\ttfamily\small,
  keywordstyle=\color{KeywordColor}\bfseries,
  commentstyle=\color{CommentColor}\itshape,
  stringstyle=\color{StringColor},
  breaklines=true,
  frame=single,
  framesep=4pt,
  rulecolor=\color{CodeBorder},
  showstringspaces=false,
  tabsize=2,
}

\lstdefinestyle{cstyle}{
  language=C,
  backgroundcolor=\color{CodeBg},
  basicstyle=\ttfamily\small,
  keywordstyle=\color{KeywordColor}\bfseries,
  commentstyle=\color{CommentColor}\itshape,
  stringstyle=\color{StringColor},
  numbers=left,
  numberstyle=\tiny\color{MidGray},
  numbersep=8pt,
  breaklines=true,
  frame=single,
  framesep=4pt,
  rulecolor=\color{CodeBorder},
  showstringspaces=false,
  tabsize=4,
}

\lstdefinestyle{yamlstyle}{
  language=,
  backgroundcolor=\color{LightGray},
  basicstyle=\ttfamily\small,
  keywordstyle=\color{KeywordColor}\bfseries,
  commentstyle=\color{CommentColor}\itshape,
  stringstyle=\color{StringColor},
  breaklines=true,
  frame=single,
  framesep=4pt,
  rulecolor=\color{CodeBorder},
  showstringspaces=false,
}

\lstset{style=pythonstyle}

% ─── Custom boxes ────────────────────────────────────────────
\tcbset{
  notestyle/.style={
    enhanced, breakable,
    colback=NoteBg, colframe=NoteBorder,
    boxrule=1pt, arc=3pt,
    fonttitle=\bfseries\small\color{white},
    coltitle=white,
    attach boxed title to top left={yshift=-2mm, xshift=5mm},
    boxed title style={colback=NoteBorder, arc=2pt},
  },
  tipstyle/.style={
    enhanced, breakable,
    colback=TipBg, colframe=TipBorder,
    boxrule=1pt, arc=3pt,
    fonttitle=\bfseries\small,
    attach boxed title to top left={yshift=-2mm, xshift=5mm},
    boxed title style={colback=TipBorder, coltext=white, arc=2pt},
  },
  warnstyle/.style={
    enhanced, breakable,
    colback=WarnBg, colframe=WarnBorder,
    boxrule=1pt, arc=3pt,
    fonttitle=\bfseries\small,
    attach boxed title to top left={yshift=-2mm, xshift=5mm},
    boxed title style={colback=WarnBorder, coltext=white, arc=2pt},
  },
  analogystyle/.style={
    enhanced, breakable,
    colback={rgb,255:red,250;green,240;blue,255},
    colframe={rgb,255:red,120;green,80;blue,180},
    boxrule=1pt, arc=4pt,
    fonttitle=\bfseries\small\itshape,
    attach boxed title to top left={yshift=-2mm, xshift=5mm},
    boxed title style={
      colback={rgb,255:red,120;green,80;blue,180},
      coltext=white, arc=2pt},
  },
}

\newcommand{\notebox}[2]{\begin{tcolorbox}[notestyle, title={#1}]#2\end{tcolorbox}}
\newcommand{\tipbox}[2]{\begin{tcolorbox}[tipstyle, title={#1}]#2\end{tcolorbox}}
\newcommand{\warnbox}[2]{\begin{tcolorbox}[warnstyle, title={#1}]#2\end{tcolorbox}}
\newcommand{\analogybox}[2]{\begin{tcolorbox}[analogystyle, title={Analogy: #1}]#2\end{tcolorbox}}

% ─── Learning objectives box ─────────────────────────────────
\newenvironment{learningobjectives}
  {\begin{tcolorbox}[enhanced, breakable,
    colback={rgb,255:red,232;green,240;blue,255},
    colframe=NavyBlue, boxrule=1.5pt, arc=4pt,
    title={\bfseries Learning Objectives},
    fonttitle=\bfseries\color{white},
    attach boxed title to top left={yshift=-2mm, xshift=5mm},
    boxed title style={colback=NavyBlue, arc=2pt}]
  \begin{itemize}[leftmargin=1.5em, itemsep=2pt]}
  {\end{itemize}\end{tcolorbox}}

% ─── Summary box ─────────────────────────────────────────────
\newenvironment{chapsummary}
  {\begin{tcolorbox}[enhanced, breakable,
    colback={rgb,255:red,245;green,245;blue,225},
    colframe=AccentOrange, boxrule=1.5pt, arc=4pt,
    title={\bfseries Chapter Summary},
    fonttitle=\bfseries\color{white},
    attach boxed title to top left={yshift=-2mm, xshift=5mm},
    boxed title style={colback=AccentOrange, arc=2pt}]}
  {\end{tcolorbox}}

% ─── Exercise box ─────────────────────────────────────────────
\newcounter{exercise}[chapter]
\newenvironment{exercise}[1][]
  {\refstepcounter{exercise}
   \begin{tcolorbox}[enhanced, breakable,
    colback={rgb,255:red,248;green,248;blue,248},
    colframe=DarkGray, boxrule=0.8pt, arc=3pt,
    title={\bfseries Exercise \thechapter.\theexercise\ifx&#1&\else: #1\fi},
    fonttitle=\bfseries\small,
    attach boxed title to top left={yshift=-2mm, xshift=5mm},
    boxed title style={colback=DarkGray, coltext=white, arc=2pt}]}
  {\end{tcolorbox}}

% ─── Key term ────────────────────────────────────────────────
\newcommand{\keyterm}[1]{\textbf{\color{AccentOrange}#1}}
\newcommand{\code}[1]{\texttt{#1}}

% ─── Part title style ────────────────────────────────────────
\titleformat{\part}[display]
  {\normalfont\Huge\bfseries\centering\color{NavyBlue}}
  {\color{AccentOrange}\MakeUppercase{Part}\ \thepart}
  {1em}
  {}

% ─── Misc ────────────────────────────────────────────────────
\setlength{\parskip}{4pt}
\setlength{\parindent}{1.5em}
\captionsetup{font=small, labelfont={bf,color=NavyBlue}}

% ============================================================
%  DOCUMENT START
% ============================================================
\begin{document}

% ─── Half-title ─────────────────────────────────────────────
\thispagestyle{empty}
\vspace*{4cm}
\begin{center}
  {\Huge\bfseries\color{NavyBlue} AI Inference Engineering}\\[0.6em]
  {\Large\color{AccentOrange} Building Fast, Scalable, and\\[0.2em]
   Production-Ready AI Systems}
\end{center}
\clearpage

% ─── Title page ─────────────────────────────────────────────
\thispagestyle{empty}
\begin{tikzpicture}[remember picture, overlay]
  \fill[NavyBlue] (current page.north west) rectangle
    ([yshift=-5cm]current page.north east);
  \node[anchor=north, text=white, font=\LARGE\bfseries, inner sep=0pt]
    at ([yshift=-1.5cm]current page.north) {AI Inference Engineering};
  \node[anchor=north, text=AccentOrange, font=\large, inner sep=0pt]
    at ([yshift=-2.8cm]current page.north)
    {Building Fast, Scalable, and Production-Ready AI Systems};
\end{tikzpicture}

\vspace*{5.5cm}

\begin{center}
  \begin{tikzpicture}[scale=0.9]
    % Simple GPU chip diagram
    \draw[fill=NavyBlue!10, draw=NavyBlue, line width=2pt, rounded corners=6pt]
      (0,0) rectangle (8,5);
    \foreach \x in {0.5,1.5,2.5,3.5,4.5,5.5,6.5,7.5}
      \foreach \y in {0.5,1.5,2.5,3.5,4.5}
        \draw[fill=NavyBlue!60, draw=NavyBlue, rounded corners=2pt]
          (\x-0.38,\y-0.38) rectangle (\x+0.38,\y+0.38);
    \node[text=white, font=\large\bfseries] at (4,2.5) {GPU};
    % arrows into GPU
    \draw[-{Stealth[length=6pt]}, AccentOrange, line width=2pt]
      (-1.5,2.5) -- (0,2.5) node[midway,above,font=\small,color=AccentOrange]
      {Requests};
    \draw[-{Stealth[length=6pt]}, AccentGreen, line width=2pt]
      (8,2.5) -- (9.5,2.5) node[midway,above,font=\small,color=AccentGreen]
      {Tokens};
  \end{tikzpicture}
\end{center}

\vspace{1cm}
\begin{center}
  {\Large\color{NavyBlue} From First Principles to Production Systems}\\[1em]
  {\normalsize\color{MidGray}
   Covers: GPU Architecture $\cdot$ CUDA $\cdot$ Transformers $\cdot$ KV Cache $\cdot$
   Quantization\\
   vLLM $\cdot$ TensorRT-LLM $\cdot$ SGLang $\cdot$ Speculative Decoding $\cdot$
   Distributed Inference\\
   Kubernetes $\cdot$ Observability $\cdot$ NVIDIA H100/Blackwell $\cdot$ 2025--2026
  }
\end{center}

\clearpage

% ─── Copyright page ─────────────────────────────────────────
\thispagestyle{empty}
\vspace*{\fill}
{\small
\noindent\textbf{AI Inference Engineering: Building Fast, Scalable, and Production-Ready AI Systems}\\[0.5em]
\noindent Edition 1.0 --- May 2026\\[0.5em]
\noindent This book covers GPU inference engineering from first principles through
production deployment.  It incorporates the latest developments in NVIDIA Blackwell
architecture (B200/GB200), vLLM PagedAttention, TensorRT-LLM v1.2, SGLang, speculative
decoding, FP8/FP4 quantization, and NVIDIA Dynamo disaggregated serving.\\[0.5em]
\noindent \textit{For educational purposes.}\\[1em]
\noindent \textbf{Typeset with} \LaTeX{} / TeX Live 2023.
}
\vspace*{\fill}
\clearpage

% ─── Table of contents ───────────────────────────────────────
\tableofcontents
\clearpage

% ─── Preface ─────────────────────────────────────────────────
\chapter*{Preface}
\addcontentsline{toc}{chapter}{Preface}
\markboth{Preface}{Preface}

\noindent\textit{"The best way to learn is to teach."}

\vspace{0.5em}

This book exists because AI is no longer just a research topic --- it is an
engineering discipline.  Every time you ask a question to ChatGPT, Claude, or
Gemini, thousands of GPU calculations happen in milliseconds.  Behind that seamless
experience is a carefully engineered system: one that balances speed, cost, memory,
and reliability at a massive scale.

\emph{AI Inference Engineering} teaches you how to build and operate those systems.

\bigskip
\noindent\textbf{Who this book is for}

\begin{itemize}
  \item Software engineers wanting to understand AI infrastructure
  \item ML engineers who want to go beyond model training
  \item Platform engineers deploying AI at scale
  \item Students learning GPU systems and distributed computing
  \item Anyone curious about what happens after a model is trained
\end{itemize}

\bigskip
\noindent\textbf{What you will learn}

You will journey from the very first question --- \textit{what is a GPU?} --- all the
way to operating a multi-node, multi-GPU inference cluster serving millions of
requests per day.  Every concept is built from scratch using simple language and
real-world analogies before advancing to production-grade engineering detail.

\bigskip
\noindent\textbf{How to read this book}

Read it linearly the first time.  Each chapter builds on the previous one.
Revisit individual chapters as reference when working on production systems.
Run the code examples --- intuition comes from doing.

\bigskip
\noindent\textbf{A note on versions}

The field moves fast.  This edition reflects the state of the art as of mid-2026:
NVIDIA H100 and Blackwell (B200/GB200) GPUs, vLLM v0.7+, TensorRT-LLM v1.2,
SGLang v0.4+, and Kubernetes-based deployment patterns.  Core concepts --- batching,
memory management, quantization, parallelism --- are timeless.

\clearpage

% ============================================================
%  PART I: FOUNDATIONS
% ============================================================
\part{Foundations}

% ============================================================
\chapter{Why Inference Engineering Matters}
\label{ch:why}
% ============================================================

\begin{learningobjectives}
  \item Understand the difference between AI training and AI inference
  \item Learn why inference efficiency is the key challenge for production AI
  \item Grasp the scale of modern AI inference systems
  \item Appreciate the engineering disciplines that make AI products work
\end{learningobjectives}

% ─────────────────────────────────────────────────────────────
\section{The Two Lives of an AI Model}

Every AI model lives two lives.

In its \textbf{first life}, the model learns.  Enormous clusters of GPUs process
billions of text samples, adjusting billions of numerical parameters over weeks or
months.  This is called \keyterm{training}\index{training}.

In its \textbf{second life}, the model works.  It answers questions, writes code,
summarizes documents, and assists doctors, lawyers, and students around the world.
This is called \keyterm{inference}\index{inference}.

Training is a one-time event.  Inference happens billions of times a day.

\analogybox{The Chef and the Restaurant}{
Imagine a chef who spends years learning to cook --- practising recipes, mastering
techniques, developing a personal style.  That long learning period is
\textit{training}.

Now the restaurant opens.  Every day the chef must cook hundreds of meals, each one
served in minutes.  That daily cooking is \textit{inference}.

The chef trains once.  The restaurant runs forever.
}

% ─────────────────────────────────────────────────────────────
\section{The Inference Problem}

When an AI model answers a question, it performs billions of mathematical
operations for every sentence it generates.  A single forward pass through a
70-billion-parameter model involves roughly 140 billion floating-point
multiplications.

\bigskip
\noindent Now multiply that by the users:

\begin{itemize}
  \item ChatGPT processes over 100 million queries per day.
  \item A single query might generate 200--1000 tokens.
  \item Each token requires a full model forward pass.
\end{itemize}

\noindent This means AI companies run \textbf{trillions} of mathematical operations
every second just to keep their services alive.  If those operations are not
performed efficiently, one of two things happens: the service becomes unbearably slow,
or the cost becomes unbearably high.

\notebox{Key Insight}{
NVIDIA's 2025 Annual Report explicitly noted that ``inference workloads surpassed
training'' in revenue terms.  The industry has shifted: getting models to run
efficiently is now more important commercially than making them more accurate.
}

% ─────────────────────────────────────────────────────────────
\section{What Inference Engineers Do}

An inference engineer's job is to make AI fast, cheap, and reliable at scale.
This involves:

\begin{itemize}
  \item \textbf{Hardware selection:} Which GPUs to buy or rent, and how many
  \item \textbf{Runtime optimization:} Making each GPU operation as fast as possible
  \item \textbf{Memory management:} Fitting large models into limited GPU memory
  \item \textbf{Batching:} Serving many users simultaneously from one GPU
  \item \textbf{Quantization:} Compressing model weights to use less memory and run faster
  \item \textbf{Deployment:} Packaging and scaling inference services on Kubernetes
  \item \textbf{Observability:} Monitoring latency, throughput, and cost in production
\end{itemize}

% ─────────────────────────────────────────────────────────────
\section{The Three Metrics That Matter}

\index{latency}\index{throughput}\index{memory}
Every inference system is measured by three core metrics:

\begin{description}
  \item[\keyterm{Latency}] How long does the user wait?  Specifically,
    \textit{Time to First Token (TTFT)}\index{TTFT} --- the delay before the
    first word appears --- and \textit{Time per Output Token (TPOT)}\index{TPOT}
    --- how quickly subsequent words arrive.
  \item[\keyterm{Throughput}] How many requests can the system process per second?
    Higher throughput means lower cost per query.
  \item[\keyterm{Memory}] How much GPU memory does the system use?  Memory limits
    how large a model you can serve and how many requests you can handle
    simultaneously.
\end{description}

These three metrics form a \textbf{triangle of tradeoffs}.  Improving one often
hurts another.  A large batch size improves throughput but increases latency.
Quantization reduces memory but may reduce accuracy.  Understanding these tradeoffs
is the core skill of an inference engineer.

% ─────────────────────────────────────────────────────────────
\section{A Preview of the Journey}

This book walks through the entire inference engineering stack, layer by layer:

\begin{center}
\begin{tikzpicture}[node distance=0.5cm]
  \tikzstyle{layer}=[rectangle, draw=NavyBlue, fill=NavyBlue!8,
                     rounded corners=4pt, minimum width=7cm,
                     minimum height=0.7cm, font=\small]
  \node[layer] (app) {Chapter 10: Production Deployment \& Kubernetes};
  \node[layer, below=of app] (obs) {Chapter 9: Observability \& Profiling};
  \node[layer, below=of obs] (dist) {Chapter 8: Distributed Inference};
  \node[layer, below=of dist] (rt) {Chapter 7: Inference Runtimes (vLLM, TRT-LLM, SGLang)};
  \node[layer, below=of rt] (sd) {Chapter 6: Speculative Decoding};
  \node[layer, below=of sd] (quant) {Chapter 5: Quantization};
  \node[layer, below=of quant] (kv) {Chapter 4: KV Cache \& Memory Management};
  \node[layer, below=of kv] (tx) {Chapter 3: Transformer Inference Mechanics};
  \node[layer, below=of tx] (cuda) {Chapter 2: GPU Architecture \& CUDA};
  \node[layer, below=of cuda, fill=AccentOrange!15, draw=AccentOrange]
    (found) {Chapter 1: Why Inference Engineering Matters (this chapter)};
\end{tikzpicture}
\end{center}

Each layer depends on the one below it.  Build a firm foundation and the rest
becomes intuitive.

\begin{chapsummary}
AI inference is the act of running a trained model to answer user queries.
It is performed billions of times daily and requires careful engineering to be fast
and affordable.  Three key metrics --- latency, throughput, and memory --- govern
every decision an inference engineer makes.  This book builds your knowledge from
the GPU chip upward to a full production deployment.
\end{chapsummary}

\begin{exercise}[Reflection]
Before reading further, write down three AI products you use daily.  For each one,
estimate: How many users might use it simultaneously?  What latency would feel
``good enough'' to you?  Keep these numbers in mind as you read --- they will help
you appreciate why every optimization in this book matters.
\end{exercise}

% ============================================================
\chapter{GPU Architecture and CUDA Fundamentals}
\label{ch:gpu}
% ============================================================

\begin{learningobjectives}
  \item Understand the fundamental difference between CPUs and GPUs
  \item Learn GPU memory hierarchy: registers, shared memory, HBM
  \item Understand CUDA programming model: threads, blocks, grids
  \item Grasp how Tensor Cores accelerate matrix operations
  \item Understand NVIDIA's GPU generations: Volta through Blackwell
\end{learningobjectives}

% ─────────────────────────────────────────────────────────────
\section{Why GPUs? The Parallel Computation Story}

A CPU\index{CPU} (Central Processing Unit) is designed for sequential, versatile
work.  It has a small number of very powerful cores --- typically 8 to 64 ---
each capable of complex decision-making, branch prediction, and running operating
systems.

A GPU\index{GPU} (Graphics Processing Unit) is built for a completely different
purpose: doing the \textit{same simple operation thousands of times simultaneously}.

\analogybox{One Expert vs.\ One Thousand Students}{
Imagine you need to add the numbers on 10{,}000 pieces of paper.

A \textbf{CPU} is like one expert accountant.  Fast, smart, capable of handling
complex calculations.  But it works through papers one at a time --- or maybe a
few at a time with multiple cores.

A \textbf{GPU} is like 10{,}000 students with calculators, each working on one
paper simultaneously.  Each student is slower and less capable than the expert.
But together, they finish the whole pile in the time the expert handles one.

For AI, which involves performing the same multiplication millions of times across
different data, the 10{,}000 students win --- by orders of magnitude.
}

% ─────────────────────────────────────────────────────────────
\section{GPU Hardware Architecture}

\subsection{Streaming Multiprocessors (SMs)}

The basic computational unit of a GPU is the
\keyterm{Streaming Multiprocessor}\index{SM, Streaming Multiprocessor} (SM).
An NVIDIA H100 GPU has 132 SMs.  Each SM contains:

\begin{itemize}
  \item \textbf{CUDA cores:} General-purpose floating-point and integer compute units
  \item \textbf{Tensor Cores:} Specialized matrix-multiply units (described below)
  \item \textbf{Registers:} Tiny, ultra-fast memory local to each thread
  \item \textbf{Shared Memory / L1 Cache:} Fast memory shared across threads in one SM
  \item \textbf{Warp Schedulers:} Hardware that manages thread execution
\end{itemize}

\begin{center}
\begin{tikzpicture}[scale=0.85]
  % GPU chip outline
  \draw[fill=NavyBlue!8, draw=NavyBlue, line width=1.5pt, rounded corners=8pt]
    (0,0) rectangle (12,8);
  \node[font=\large\bfseries\color{NavyBlue}, anchor=north] at (6,7.8) {GPU Die};

  % GPC blocks
  \foreach \gx/\gy/\gn in {0.5/4.2/GPC 1, 4.5/4.2/GPC 2, 8.5/4.2/GPC 3,
                             0.5/0.3/GPC 4, 4.5/0.3/GPC 5, 8.5/0.3/GPC 6} {
    \draw[fill=NavyBlue!20, draw=NavyBlue, rounded corners=4pt]
      (\gx,\gy) rectangle (\gx+3.5,\gy+3.5);
    \node[font=\tiny\bfseries\color{NavyBlue}, anchor=north] at (\gx+1.75,\gy+3.4) {\gn};
    % SM inside GPC
    \foreach \sx in {0,1} \foreach \sy in {0,1,2} {
      \draw[fill=AccentOrange!30, draw=AccentOrange, rounded corners=2pt]
        (\gx+0.2+\sx*1.65, \gy+0.25+\sy*1.0) rectangle
        (\gx+1.55+\sx*1.65, \gy+1.05+\sy*1.0);
      \node[font=\tiny, anchor=center] at
        (\gx+0.875+\sx*1.65, \gy+0.65+\sy*1.0) {SM};
    }
  }

  % HBM label
  \draw[fill=AccentGreen!20, draw=AccentGreen, line width=1pt]
    (-1.8, 1) rectangle (-0.2, 7);
  \node[font=\small\bfseries\color{AccentGreen}, rotate=90] at (-1,4) {HBM3e};
  \draw[-{Stealth}, AccentGreen, line width=1pt] (-0.2,4) -- (0,4);

  \node[font=\tiny\color{MidGray}, anchor=south] at (6,-0.1)
    {Simplified: H100 has 132 SMs across 8 GPCs};
\end{tikzpicture}
\end{center}

\subsection{The GPU Memory Hierarchy}

\index{memory hierarchy}
Memory in a GPU is organized in layers, from fastest (but tiny) to slowest
(but huge):

\begin{longtable}{@{}p{3cm}p{2cm}p{2cm}p{5.5cm}@{}}
  \toprule
  \textbf{Memory Type} & \textbf{Size} & \textbf{Bandwidth} & \textbf{Purpose} \\
  \midrule
  \endfirsthead
  \toprule
  \textbf{Memory Type} & \textbf{Size} & \textbf{Bandwidth} & \textbf{Purpose} \\
  \midrule
  \endhead
  Registers & 256 KB/SM & $\sim$20 TB/s & Per-thread data \\
  L1/Shared Memory & 228 KB/SM & $\sim$19 TB/s & Thread cooperation within SM \\
  L2 Cache & 50 MB & $\sim$12 TB/s & Cross-SM caching \\
  HBM3e (VRAM)\index{HBM} & 80--141 GB & 3.35 TB/s & Model weights, KV cache \\
  CPU RAM (via NVLink/PCIe) & TBs & 900 GB/s & Offload, overflow \\
  \bottomrule
  \caption{H100 GPU Memory Hierarchy}
\end{longtable}

\tipbox{The Golden Rule of GPU Programming}{
The fastest code is code that keeps data as close to the compute units as possible.
If your CUDA kernel reads from HBM when it could read from shared memory, you are
leaving 5$\times$ performance on the table.  Every serious GPU optimization is,
at its heart, a memory access optimization.
}

\subsection{Tensor Cores: The Matrix Multiplication Engines}

\keyterm{Tensor Cores}\index{Tensor Cores} are specialized hardware units inside
each SM designed for one specific task: ultra-fast matrix multiplication.  They
were introduced in the Volta architecture (2017) and have been enhanced in every
generation since.

Why does matrix multiplication matter for AI?  Because the core operation in every
neural network layer is a \textbf{matrix multiplication} (GEMM --- General Matrix
Multiply).  When you multiply an input vector by a weight matrix to produce an
output, that is a matrix multiply.  Modern LLMs perform this operation hundreds of
times per forward pass.

\begin{center}
\begin{tikzpicture}[scale=0.9]
  % Matrix A
  \draw[fill=NavyBlue!15, draw=NavyBlue] (0,0) rectangle (1.5,2);
  \node[font=\small\bfseries] at (0.75,1) {A};
  \node[font=\tiny\color{MidGray}] at (0.75,-0.3) {$m \times k$};

  % times
  \node[font=\large] at (2.0,1) {$\times$};

  % Matrix B
  \draw[fill=AccentOrange!20, draw=AccentOrange] (2.5,0.5) rectangle (4.5,2.5);
  \node[font=\small\bfseries] at (3.5,1.5) {B};
  \node[font=\tiny\color{MidGray}] at (3.5,-0.3) {$k \times n$};

  % equals
  \node[font=\large] at (5.0,1) {$=$};

  % Matrix C
  \draw[fill=AccentGreen!20, draw=AccentGreen] (5.5,0) rectangle (7.5,2);
  \node[font=\small\bfseries] at (6.5,1) {C};
  \node[font=\tiny\color{MidGray}] at (6.5,-0.3) {$m \times n$};

  % Tensor Core label
  \draw[-{Stealth}, thick, AccentOrange] (8,1.5) -- (6.8,1.5);
  \node[anchor=west, font=\small\color{AccentOrange}] at (8,1.5)
    {Tensor Core computes this};
  \node[anchor=west, font=\tiny\color{MidGray}] at (8,1.0)
    {in a single hardware cycle};
\end{tikzpicture}
\end{center}

The H100 Tensor Core operates on \textbf{FP8} precision natively, performing a
$16 \times 8 \times 16$ matrix multiplication in a single instruction --- the
equivalent of 2048 floating-point multiply-add operations per cycle.

\subsubsection{GPU Generations and Precision Support}

\begin{longtable}{@{}p{2cm}p{2cm}p{1.5cm}p{1.5cm}p{5cm}@{}}
  \toprule
  \textbf{GPU} & \textbf{Arch.} & \textbf{Year} &
  \textbf{Peak (FP8)} & \textbf{Key Features} \\
  \midrule
  \endfirsthead
  \toprule
  \textbf{GPU} & \textbf{Arch.} & \textbf{Year} &
  \textbf{Peak (FP8)} & \textbf{Key Features} \\
  \midrule
  \endhead
  A100 & Ampere & 2020 & 624 TFLOPS & First FP16 Tensor Cores, NVLink 3 \\
  H100 SXM & Hopper & 2022 & 3958 TFLOPS & FP8 native, Transformer Engine, NVLink 4 \\
  H200 SXM & Hopper+ & 2024 & 3958 TFLOPS & 141 GB HBM3e memory \\
  B200 & Blackwell & 2025 & 9000 TFLOPS & FP4, 5th-gen Tensor Cores, 30$\times$ H100 inference \\
  GB200 NVL72 & Blackwell & 2025 & $\sim$1 EFLOPS & 72 GPUs in a rack, disaggregated serving \\
  \bottomrule
  \caption{NVIDIA GPU Generations for AI Inference}
  \label{tab:gpus}
\end{longtable}
\index{H100}\index{Blackwell}\index{B200}

\notebox{Blackwell's Leap}{
The NVIDIA B200 (Blackwell) delivers up to 30$\times$ the inference performance
of the H100 on specific workloads, particularly for large models using FP4
quantization.  The fifth-generation Tensor Cores support FP32, FP16, BF16, FP8,
INT8, FP6, and FP4 --- enabling extremely flexible precision choices.
}

% ─────────────────────────────────────────────────────────────
\section{CUDA Programming Model}
\index{CUDA}

\keyterm{CUDA} (Compute Unified Device Architecture) is NVIDIA's programming
framework for writing code that runs on GPUs.  Every inference runtime ---
vLLM, TensorRT, PyTorch --- ultimately executes CUDA code under the hood.

\subsection{The Thread Hierarchy}

CUDA organizes GPU execution in a strict three-level hierarchy:

\begin{itemize}
  \item \textbf{Thread:\index{thread}} The smallest unit.  One thread performs one computation.
  \item \textbf{Block:\index{block}} A group of up to 1024 threads that share memory and can synchronize.
  \item \textbf{Grid:\index{grid}} A collection of blocks that together run one kernel.
\end{itemize}

\analogybox{School Analogy}{
Think of a grid as a \textbf{school}.
Each block is a \textbf{classroom}.
Each thread is a \textbf{student}.

Students in one classroom can pass notes (shared memory).
Students in different classrooms cannot talk directly.
The principal (GPU scheduler) tells each classroom what lesson to do (\textit{kernel}).
}

\subsection{A Simple CUDA Kernel}

\begin{lstlisting}[style=cstyle, caption={A simple CUDA vector addition kernel},
                   label={lst:cuda_add}]
// This function runs on the GPU -- one thread per element
__global__ void vector_add(float* A, float* B, float* C, int N) {
    // Each thread computes its own index
    // blockIdx.x  = which block am I in?
    // blockDim.x  = how many threads per block?
    // threadIdx.x = which thread am I within my block?
    int idx = blockIdx.x * blockDim.x + threadIdx.x;

    // Boundary check: don't go past the array end
    if (idx < N) {
        C[idx] = A[idx] + B[idx];   // Each thread adds one element
    }
}

// Called from the CPU (host)
void run_vector_add(float* A, float* B, float* C, int N) {
    int threads_per_block = 256;
    int num_blocks = (N + threads_per_block - 1) / threads_per_block;

    // <<< grid_size, block_size >>> launches the kernel
    vector_add<<<num_blocks, threads_per_block>>>(A, B, C, N);
}
\end{lstlisting}

\subsection{Memory Management in CUDA}

\begin{lstlisting}[style=cstyle, caption={CUDA memory allocation and transfer}]
// Allocate memory on CPU (host)
float* h_A = (float*)malloc(N * sizeof(float));

// Allocate memory on GPU (device)
float* d_A;
cudaMalloc(&d_A, N * sizeof(float));    // GPU memory allocation

// Transfer data CPU -> GPU
cudaMemcpy(d_A, h_A, N*sizeof(float), cudaMemcpyHostToDevice);

// Run kernel on GPU
my_kernel<<<grid, block>>>(d_A, ...);

// Transfer results GPU -> CPU
cudaMemcpy(h_result, d_result, N*sizeof(float), cudaMemcpyDeviceToHost);

// Free GPU memory
cudaFree(d_A);
\end{lstlisting}

\notebox{Memory Transfer is Expensive}{
Moving data between CPU and GPU via PCIe costs time --- typically 16 GB/s on a
PCIe 4.0 $\times$16 connection.  This is 100$\times$ slower than reading from GPU
HBM.  In production inference, the goal is to keep \emph{everything} on the GPU
and minimize CPU-GPU transfers.
}

\subsection{Warps: The Real Execution Unit}

While blocks are the logical unit, hardware executes threads in groups of 32
called \keyterm{warps}\index{warp}.  All 32 threads in a warp execute the same
instruction simultaneously.  If threads in a warp take different code paths
(e.g., an if/else), both paths execute sequentially with some threads masked off.
This is called \keyterm{warp divergence}\index{warp divergence} and is a major
source of GPU performance loss.

\begin{exercise}[Thread Index Calculation]
A CUDA kernel is launched with a grid of 16 blocks, each containing 256 threads.
\begin{enumerate}
  \item What is the total number of threads?
  \item If \texttt{blockIdx.x = 5} and \texttt{threadIdx.x = 100},
        what is the global thread index?
  \item If processing an array of 3000 elements, how many threads will be out
        of bounds and need a guard check?
\end{enumerate}
\end{exercise}

% ─────────────────────────────────────────────────────────────
\section{FlashAttention: The Kernel That Changed Inference}
\index{FlashAttention}

Before we leave GPU basics, let's preview one of the most important CUDA
optimizations ever invented for AI: \keyterm{FlashAttention}.

The standard attention computation in a Transformer reads the key, query, and
value matrices from HBM, computes a large intermediate matrix (the attention
scores), writes it back to HBM, then reads it again to apply softmax.  This is
\textbf{extremely memory-bandwidth-bound}.

FlashAttention (Dao et al., 2022) reorganizes this computation so that all
intermediate results stay in shared memory.  It never writes the large intermediate
matrix to HBM at all.  This single change achieves:

\begin{itemize}
  \item $2-4\times$ faster attention on GPU
  \item $5-20\times$ less memory for attention
  \item Exact same mathematical result as standard attention
\end{itemize}

FlashAttention is now standard in every production inference engine.

\begin{chapsummary}
GPUs accelerate AI because they perform thousands of operations in parallel.
The memory hierarchy (registers $\to$ shared memory $\to$ HBM) is central to
performance: keep data close to compute.  CUDA organizes execution as grids of
blocks of threads.  Tensor Cores handle matrix multiplication at extreme speed.
Modern H100 and Blackwell GPUs support FP8/FP4 precision for even higher
throughput.  FlashAttention shows how smart kernel design can dramatically improve
both speed and memory use.
\end{chapsummary}

% ============================================================
\chapter{Transformer Inference Mechanics}
\label{ch:transformer}
% ============================================================

\begin{learningobjectives}
  \item Understand how Transformer models generate text token by token
  \item Learn the two phases of inference: prefill and decode
  \item Understand why decode is memory-bandwidth-bound, not compute-bound
  \item Grasp the roofline model for compute vs.\ memory bound analysis
\end{learningobjectives}

% ─────────────────────────────────────────────────────────────
\section{How Transformers Generate Text}

\index{autoregressive decoding}
Large Language Models generate text \textbf{one token at a time}.  A
\keyterm{token}\index{token} is roughly a word-piece --- ``inference'' might be
one token, while ``transformers'' might be two: ``transform'' and ``ers''.

When you send a prompt to an LLM, here is what happens:

\begin{enumerate}
  \item Your prompt is tokenized into a list of integer token IDs.
  \item The model processes all prompt tokens simultaneously (\textbf{prefill phase}).
  \item The model generates one output token, then another, then another, until
        it produces an end-of-sequence token or reaches a length limit
        (\textbf{decode phase}).
\end{enumerate}

\analogybox{The Sentence Completion Game}{
Imagine you are playing a word-completion game.  You read the entire sentence so
far, then you predict the most likely next word.  You write it down.  Now you read
the whole sentence again, including the word you just wrote, and predict the next
word.  Repeat.

This is exactly how an LLM works.  The entire conversation history is reread every
time a new token is generated.  This is called \textit{autoregressive} generation.
}

\subsection{The Two Phases in Detail}

\keyterm{Prefill}\index{prefill} processes the entire input prompt in a single
forward pass.  If your prompt has 512 tokens, all 512 are processed at once,
using the GPU's full parallel compute power.  Prefill is \textbf{compute-bound}:
the GPU is busy multiplying matrices.

\keyterm{Decode}\index{decode} generates one token per step.  At each step, only a
single new token is processed (the most recently generated one), while all
previous tokens are accessible through the KV cache (Chapter~\ref{ch:kvcache}).
Decode is \textbf{memory-bandwidth-bound}: the GPU spends most of its time loading
model weights from HBM, not computing.

\begin{center}
\begin{tikzpicture}[scale=0.9]
  % Timeline
  \draw[-{Stealth}, line width=1.5pt] (0,0) -- (12,0) node[right] {Time};

  % Prefill block
  \draw[fill=NavyBlue!25, draw=NavyBlue, rounded corners=4pt]
    (0,0.2) rectangle (3,1.5);
  \node[font=\small\bfseries\color{NavyBlue}] at (1.5,0.85)
    {\textbf{PREFILL}};
  \node[font=\tiny\color{MidGray}] at (1.5,0.4)
    {All prompt tokens processed};
  \node[font=\tiny\color{MidGray}] at (1.5,0.2)
    {Compute-bound};

  % Decode blocks
  \foreach \i/\tok in {1/``The'', 2/``sky'', 3/``is'', 4/``blue'', 5/``.'', 6/\small{\textless /s\textgreater}} {
    \pgfmathtruncatemacro{\x}{3 + (\i-1)*1.4}
    \draw[fill=AccentOrange!20, draw=AccentOrange, rounded corners=3pt]
      (\x,0.2) rectangle (\x+1.2,1.5);
    \node[font=\tiny\bfseries\color{AccentOrange}] at (\x+0.6,1.0) {\tok};
    \node[font=\tiny\color{MidGray}] at (\x+0.6,0.45) {step \i};
  }

  % TTFT arrow
  \draw[{Stealth}-{Stealth}, MidGray] (0,-0.4) -- (3,-0.4);
  \node[font=\tiny\color{MidGray}] at (1.5,-0.65) {TTFT};

  % TPOT arrow
  \draw[{Stealth}-{Stealth}, MidGray] (3,-0.4) -- (4.4,-0.4);
  \node[font=\tiny\color{MidGray}] at (3.7,-0.65) {TPOT};
\end{tikzpicture}
\end{center}

% ─────────────────────────────────────────────────────────────
\section{The Roofline Model: Compute vs.\ Memory Bound}
\index{roofline model}

To understand why decode is slow, we need the \keyterm{roofline model} --- a
simple framework for determining whether a computation is limited by arithmetic
speed or memory bandwidth.

Every GPU has two fundamental limits:
\begin{itemize}
  \item \textbf{Peak FLOPS:} How many floating-point operations per second it can
        perform (H100: $\sim$1979 TFLOPS in FP16)
  \item \textbf{Peak Memory Bandwidth:} How fast it can read data from HBM
        (H100: 3.35 TB/s)
\end{itemize}

For a given computation, the \textbf{arithmetic intensity} is:
\[
  \text{Arithmetic Intensity} =
  \frac{\text{FLOP performed}}{\text{bytes read from memory}}
\]

If arithmetic intensity is \textbf{above} a threshold (called the ridge point),
the operation is \textbf{compute-bound}: your GPU is fast enough but you're
not feeding it enough data.

If arithmetic intensity is \textbf{below} the ridge point, the operation is
\textbf{memory-bandwidth-bound}: your GPU is waiting for data from memory.

For decode with a single request:
\begin{itemize}
  \item We read the entire model ($\sim$140 GB for a 70B parameter FP16 model)
  \item We perform only $\sim$280 GFLOPs of computation per token
  \item Arithmetic intensity $\approx 2$ FLOP/byte --- far below the H100's
        ridge point of $\sim$590 FLOP/byte
\end{itemize}

\warnbox{The Decode Efficiency Problem}{
For a single request, the GPU is roughly \textbf{99\% idle during decode}.
We load 140 GB of model weights to produce maybe 50 bytes of output token.
This is deeply inefficient --- and it's the core reason batching is so
important in inference engineering.
}

% ─────────────────────────────────────────────────────────────
\section{Inside a Transformer Layer}

A single Transformer layer\index{Transformer layer} contains two main components:

\begin{enumerate}
  \item \textbf{Multi-Head Self-Attention (MHA):} Each token attends to every
        other token in the sequence.
  \item \textbf{Feed-Forward Network (FFN):} A two-layer MLP applied independently
        to each token.
\end{enumerate}

\subsection{The Attention Computation}

For a sequence of $n$ tokens, the attention computation is:

\[
  \text{Attention}(Q, K, V) = \text{softmax}\!\left(\frac{QK^\top}{\sqrt{d_k}}\right) V
\]

where $Q$, $K$, $V$ are the query, key, and value matrices derived from the input
by multiplying with learned weight matrices $W_Q$, $W_K$, $W_V$.

During \textbf{decode}, we only compute $Q$ for the new token, but we need all
previous $K$ and $V$ values.  This is why we need a \textbf{KV Cache}
(Chapter~\ref{ch:kvcache}).

\subsection{Memory Footprint of a Large Model}

For a Llama-3 70B model in FP16:

\begin{center}
\begin{tabular}{lrr}
  \toprule
  \textbf{Component} & \textbf{Parameters} & \textbf{Memory (FP16)} \\
  \midrule
  Embedding layers & 0.5B & 1 GB \\
  80 Attention layers & 40B & 80 GB \\
  80 FFN layers & 28B & 56 GB \\
  Other (LN, head) & 1.5B & 3 GB \\
  \midrule
  \textbf{Total} & \textbf{70B} & \textbf{$\sim$140 GB} \\
  \bottomrule
\end{tabular}
\end{center}

This model requires \textbf{two H100 80GB GPUs} just to load the weights, before
any KV cache or batch data.

\begin{chapsummary}
Transformers generate text autoregressively --- one token at a time.  The prefill
phase processes the entire prompt in parallel (compute-bound), while the decode
phase generates one token at a time (memory-bandwidth-bound).  The roofline model
helps understand why: during decode we load all model weights for minimal
computation.  Batching multiple requests together is the key to making decode
efficient.  Large models need substantial GPU memory just for their weights.
\end{chapsummary}

% ============================================================
\chapter{KV Cache and Memory Management}
\label{ch:kvcache}
% ============================================================

\begin{learningobjectives}
  \item Understand why the KV cache exists and what it stores
  \item Learn how PagedAttention eliminates KV cache fragmentation
  \item Calculate KV cache memory requirements for production systems
  \item Understand prefix caching and its performance benefits
\end{learningobjectives}

% ─────────────────────────────────────────────────────────────
\section{Why We Need a KV Cache}

Recall the decode phase: we generate one new token at a time, but each token
requires attention over \emph{all} previous tokens.

Without caching: at step $t$, to generate token $t+1$, we would recompute
all the key and value vectors for tokens $1$ through $t$.  This means
the computation grows linearly with sequence length: generating a 1000-token
response costs $1+2+3+\ldots+1000 = 500{,}500$ attention computations instead
of just 1000.

The \keyterm{KV Cache}\index{KV cache} solves this by \textbf{storing} the key
and value matrices computed for every past token.  When generating token $t+1$,
we reuse all stored $K$ and $V$ from steps $1$ through $t$ and only compute
new $K$ and $V$ for the latest token.

\analogybox{The Meeting Notes}{
Imagine you are taking notes in a long meeting.  Instead of re-reading the entire
transcript every time you want to write the next sentence, you keep a notepad of
the key points discussed so far.

Each time you need to reference the past, you glance at your notepad --- not the
entire transcript.

The KV cache is that notepad.  It stores the important ``memory'' of past tokens so
the GPU doesn't recompute it every step.
}

% ─────────────────────────────────────────────────────────────
\section{KV Cache Memory Calculation}

The KV cache size for a single sequence is:

\[
  \text{KV cache size} = 2 \times n_\text{layers} \times n_\text{heads} \times
  d_\text{head} \times \text{seq\_len} \times \text{bytes\_per\_element}
\]

The factor 2 accounts for both keys ($K$) and values ($V$).

For Llama-3 70B (FP16, 4096 sequence length):
\[
  2 \times 80 \times 64 \times 128 \times 4096 \times 2 \approx \mathbf{10.7\text{ GB}}
\]

\warnbox{Memory Shock}{
A single conversation of 4096 tokens with Llama-3 70B uses \textbf{10.7 GB} of KV
cache --- on top of the 140 GB for weights.  Serving 10 simultaneous users requires
107 GB of KV cache alone.  Efficient KV cache management is one of the biggest
challenges in production inference.
}

% ─────────────────────────────────────────────────────────────
\section{The Problem: KV Cache Fragmentation}

Before PagedAttention, every inference engine allocated KV cache memory using
a single large contiguous block per sequence.

Imagine you are running a hotel.  Every guest (request) is allocated a room block:
a contiguous range of rooms.  But you don't know in advance how long each guest
will stay.  You must reserve the maximum possible stay (max sequence length) for
each guest upfront.  The result: many rooms sit empty, but no new guests can be
booked because there's no single contiguous block large enough.

In production, this \textbf{fragmentation} means that even when 30\% of GPU
memory is technically free, it cannot be used because it is scattered in small
unusable pieces.  Studies found that before PagedAttention, 60--80\% of GPU
memory allocated for KV caches was wasted.

% ─────────────────────────────────────────────────────────────
\section{PagedAttention: Solving Fragmentation}
\index{PagedAttention}

\keyterm{PagedAttention} (Kwon et al., SOSP 2023 --- the core innovation behind
vLLM) borrows an idea from operating system virtual memory: \textbf{paging}.

Instead of one contiguous block per request, memory is divided into fixed-size
\textbf{blocks} (typically 16 or 32 tokens per block).  A request uses a
\textbf{block table} --- a list of pointers to the physical blocks it owns.
Blocks do not need to be contiguous in memory.

\begin{center}
\begin{tikzpicture}[scale=0.85]
  % Physical GPU memory
  \draw[fill=LightGray, draw=DarkGray, rounded corners=4pt]
    (0,0) rectangle (10,2.5);
  \node[font=\small\bfseries\color{DarkGray}, anchor=north] at (5,2.5)
    {Physical GPU Memory (KV Cache Pool)};

  % Blocks
  \foreach \i/\col/\lbl in {
    0/NavyBlue!40/Req A,
    1/AccentOrange!40/Req B,
    2/AccentGreen!40/Req C,
    3/NavyBlue!40/Req A,
    4/AccentOrange!40/Req B,
    5/MidGray!30/Free,
    6/AccentGreen!40/Req C,
    7/AccentGreen!40/Req C,
    8/NavyBlue!40/Req A,
    9/MidGray!30/Free}
  {
    \pgfmathsetmacro{\xpos}{\i}
    \draw[fill=\col, draw=DarkGray, rounded corners=2pt]
      (\xpos*1+0.05,0.3) rectangle (\xpos*1+0.9,2.2);
    \node[font=\tiny, rotate=90] at (\xpos*1+0.475,1.25) {\lbl};
    \node[font=\tiny\color{MidGray}, anchor=north] at (\xpos*1+0.475,0.25)
      {\i};
  }

  % Block tables
  \foreach \req/\col/\y/\blocks in {
    A/NavyBlue/4.0/{0,3,8},
    B/AccentOrange/3.0/{1,4},
    C/AccentGreen/2.0/{2,6,7}}
  {
    \node[font=\small\bfseries\color{\col}, anchor=east] at (-0.1,\y) {Req \req};
    \draw[fill=\col!15, draw=\col, rounded corners=2pt]
      (0,\y-0.3) rectangle (2.5,\y+0.3);
    \node[font=\small\color{\col}] at (1.25,\y) {blocks: \blocks};
  }

  % arrows from block table to physical memory
  \draw[-{Stealth}, NavyBlue!60, dashed] (2.5,4.0) -- (0.475,2.2);
  \draw[-{Stealth}, NavyBlue!60, dashed] (2.5,4.0) -- (3.475,2.2);
  \draw[-{Stealth}, NavyBlue!60, dashed] (2.5,4.0) -- (8.475,2.2);
\end{tikzpicture}
\end{center}

\notebox{PagedAttention Impact}{
PagedAttention reduces KV cache waste from 60--80\% down to under 5\%.
vLLM achieves 14--24$\times$ higher throughput than naive HuggingFace Transformers
and 2.2--3.5$\times$ higher than early Text Generation Inference (TGI) engines
for LLaMA models on NVIDIA GPUs.
}

% ─────────────────────────────────────────────────────────────
\section{Prefix Caching}
\index{prefix caching}

Many queries share common prefixes: the same system prompt, the same document
being analyzed, the same few-shot examples.

\keyterm{Prefix caching} (also called \textit{prompt caching} or \textit{RadixAttention}
in SGLang) stores the KV cache of common prefixes and \textbf{reuses} them across
multiple requests.

If 1000 users all send queries starting with the same 2000-token system prompt,
prefix caching means the KV computation for those 2000 tokens is performed
\textit{once}, not 1000 times.

\begin{lstlisting}[language=Python, caption={Enabling prefix caching in vLLM}]
from vllm import LLM, SamplingParams

# Initialize vLLM with prefix caching enabled
llm = LLM(
    model="meta-llama/Llama-3-70b-instruct",
    enable_prefix_caching=True,   # Enable prefix caching
    gpu_memory_utilization=0.90,  # Use 90% of GPU memory for KV pool
    tensor_parallel_size=2,       # Across 2 GPUs
)

# Many requests sharing the same system prompt will be served
# faster because the KV cache for the system prompt is reused
system_prompt = "You are a helpful coding assistant. " * 200  # long prompt

responses = llm.generate(
    [system_prompt + user_q for user_q in user_questions],
    SamplingParams(temperature=0.7, max_tokens=512)
)
\end{lstlisting}

% ─────────────────────────────────────────────────────────────
\section{KV Cache Quantization}
\index{KV cache quantization}

Since KV cache can consume more memory than the model weights themselves for long
sequences, quantizing the KV cache is a powerful optimization.

Recent research shows that INT8 quantization of the KV cache achieves a
\textbf{4$\times$ memory reduction} with reconstruction error below 0.004 and
attention score error below 0.1 even for 8K-dimensional heads.  FP8 KV
quantization is now natively supported in vLLM and TensorRT-LLM.

\begin{lstlisting}[language=Python, caption={FP8 KV cache quantization in vLLM}]
from vllm import LLM

llm = LLM(
    model="meta-llama/Llama-3-70b-instruct",
    kv_cache_dtype="fp8",         # Use FP8 for KV cache
    enable_prefix_caching=True,
    gpu_memory_utilization=0.92,
)
# FP8 KV cache reduces memory by ~2x vs FP16
# Minimal accuracy degradation in practice
\end{lstlisting}

\begin{chapsummary}
The KV cache stores computed keys and values for all past tokens, allowing
efficient autoregressive generation without recomputation.  KV cache memory
grows linearly with sequence length and can easily exceed model weight memory.
PagedAttention solves fragmentation by using block-based allocation similar to
OS virtual memory, enabling near-zero waste.  Prefix caching reuses KV state
for shared prompt prefixes.  FP8/INT8 quantization of the KV cache reduces
memory by 2--4$\times$ with minimal accuracy impact.
\end{chapsummary}

% ============================================================
\chapter{Quantization: Fitting More Into Less Memory}
\label{ch:quantization}
% ============================================================

\begin{learningobjectives}
  \item Understand what quantization is and why it matters for inference
  \item Learn the difference between INT8, FP8, INT4, FP4 quantization
  \item Understand post-training quantization (PTQ) vs.\ quantization-aware training
  \item Apply AWQ and GPTQ quantization in practice
  \item Understand the accuracy--speed--memory tradeoffs
\end{learningobjectives}

% ─────────────────────────────────────────────────────────────
\section{What Is Quantization?}

Neural network weights are typically stored in full precision:
\textbf{FP32} (32 bits per number) or \textbf{FP16/BF16} (16 bits per number).

\keyterm{Quantization}\index{quantization} replaces these high-precision numbers
with lower-precision representations --- fewer bits per number --- that are
\textbf{good enough} for inference.

\analogybox{The Photo Compression Analogy}{
Imagine a photograph stored at full resolution: 48 megapixels, 24 bits per pixel.
The file is huge --- say, 200 MB.

Now you apply JPEG compression.  The image now takes 2 MB.  When you look at it
on your phone screen, you cannot tell the difference.

Quantization is compression for AI model weights.  You reduce the number of bits
per weight.  The model takes less memory, runs faster, and (if done carefully)
produces answers that are almost identical to the original.
}

% ─────────────────────────────────────────────────────────────
\section{Number Formats: A Visual Guide}

\begin{center}
\begin{tikzpicture}[scale=0.9]
  % FP32
  \draw[fill=NavyBlue!20] (0,3) rectangle (8,3.6);
  \node[font=\tiny, anchor=west] at (0.1,3.3) {S};
  \foreach \i in {1,...,8} \draw[NavyBlue] (\i*0.2,3) -- (\i*0.2,3.6);
  \foreach \i in {1,...,7} \draw[MidGray] (1.6+\i*0.288,3) -- (1.6+\i*0.288,3.6);
  \draw[fill=NavyBlue!8] (1.6,3) rectangle (8,3.6);
  \node[font=\tiny\color{MidGray}] at (4,3.8) {FP32: 1 sign + 8 exponent + 23 mantissa = 32 bits};
  \node[font=\tiny\color{NavyBlue}, anchor=west] at (0.1,2.75) {Sign};
  \node[font=\tiny\color{NavyBlue}] at (1.0,2.75) {Exponent (8)};
  \node[font=\tiny\color{MidGray}] at (4.8,2.75) {Mantissa (23)};

  % FP16
  \draw[fill=AccentOrange!20] (0,1.8) rectangle (4,2.4);
  \node[font=\tiny\color{MidGray}] at (2,2.6)
    {FP16: 1+5+10 = 16 bits};

  % INT8
  \draw[fill=AccentGreen!20] (0,0.6) rectangle (2,1.2);
  \node[font=\tiny\color{MidGray}] at (1,1.4) {INT8: 8 bits (-128 to 127)};

  % INT4
  \definecolor{INT4Color}{RGB}{180,100,0}
  \draw[fill=INT4Color!20] (0,-0.6) rectangle (1,-0.0);
  \node[font=\tiny\color{MidGray}] at (0.5,-0.8) {INT4: 4 bits};

  % Size comparison
  \node[font=\small\bfseries\color{DarkGray}, anchor=west] at (8.5,3.3) {32 bits};
  \node[font=\small\bfseries\color{DarkGray}, anchor=west] at (8.5,2.1) {16 bits (2$\times$ smaller)};
  \node[font=\small\bfseries\color{DarkGray}, anchor=west] at (8.5,0.9) {8 bits (4$\times$ smaller)};
  \node[font=\small\bfseries\color{DarkGray}, anchor=west] at (8.5,-0.3) {4 bits (8$\times$ smaller)};
\end{tikzpicture}
\end{center}

\begin{longtable}{@{}p{1.5cm}p{1.5cm}p{1.5cm}p{4cm}p{3cm}@{}}
  \toprule
  \textbf{Format} & \textbf{Bits} & \textbf{Relative Size} &
  \textbf{Dynamic Range} & \textbf{Use Case} \\
  \midrule
  \endfirsthead
  \toprule
  \textbf{Format} & \textbf{Bits} & \textbf{Relative Size} &
  \textbf{Dynamic Range} & \textbf{Use Case} \\
  \midrule
  \endhead
  FP32 & 32 & 1$\times$ & $1.2\times10^{-38}$ to $3.4\times10^{38}$ & Training reference \\
  BF16 & 16 & 2$\times$ & Same range as FP32 & Default training/inference \\
  FP16 & 16 & 2$\times$ & $6\times10^{-5}$ to $65504$ & Common inference \\
  FP8 (E4M3) & 8 & 4$\times$ & $-448$ to $448$ & H100/Blackwell inference \\
  INT8 & 8 & 4$\times$ & $-128$ to $127$ & Weight/activation quantization \\
  INT4/FP4 & 4 & 8$\times$ & $-8$ to $7$ & Aggressive compression \\
  \bottomrule
  \caption{Number Formats for AI Inference}
\end{longtable}

% ─────────────────────────────────────────────────────────────
\section{Post-Training Quantization Methods}
\index{PTQ, post-training quantization}

\keyterm{Post-Training Quantization} (PTQ) applies quantization to an already
trained model, without any additional training.

\subsection{GPTQ: Round-to-Nearest with Correction}
\index{GPTQ}

GPTQ (Frantar et al., 2022) is a popular INT4 quantization method that:
\begin{enumerate}
  \item Quantizes weights layer by layer, in order
  \item After quantizing each weight, adjusts remaining weights to compensate
        for the error introduced
  \item Uses a small calibration dataset to guide quantization
\end{enumerate}

\begin{lstlisting}[language=Python, caption={Loading a GPTQ quantized model}]
from transformers import AutoModelForCausalLM, AutoTokenizer
from optimum.gptq import GPTQQuantizer

# Load a pre-quantized GPTQ model (4-bit)
# Models are available on HuggingFace Hub with "-GPTQ" suffix
model = AutoModelForCausalLM.from_pretrained(
    "TheBloke/Llama-2-70B-GPTQ",  # Pre-quantized 4-bit model
    device_map="auto",            # Auto-distribute across GPUs
    torch_dtype="auto",
)
# This 70B model now fits in ~35GB instead of ~140GB
\end{lstlisting}

\subsection{AWQ: Activation-Aware Weight Quantization}
\index{AWQ}

AWQ (Lin et al., 2023) is a smarter quantization technique that observes which
weights are most \textbf{important} for accuracy (by analyzing activations) and
protects them from quantization error.

Key insight: \textit{not all weights are equally important}.  About 1\% of
weights are ``salient'' --- they have large activations and dominate the model's
behavior.  AWQ keeps these weights at higher precision while aggressively
quantizing the rest.

\begin{lstlisting}[language=Python, caption={Using vLLM with AWQ quantized model}]
from vllm import LLM, SamplingParams

# Load AWQ 4-bit quantized model
llm = LLM(
    model="casperhansen/llama-3-70b-instruct-awq",
    quantization="awq",           # Tell vLLM this is an AWQ model
    tensor_parallel_size=2,       # Still need 2 GPUs for 4-bit 70B
    gpu_memory_utilization=0.92,
)

output = llm.generate(
    ["Explain transformer attention in one sentence."],
    SamplingParams(temperature=0.7, max_tokens=100)
)
print(output[0].outputs[0].text)
\end{lstlisting}

\subsection{FP8 Quantization on H100/Blackwell}
\index{FP8}

The H100 introduced native FP8 hardware support through its Transformer Engine.
FP8 requires minimal post-processing compared to INT quantization because it
maintains a floating-point representation, preserving the dynamic range needed
by neural network activations.

\begin{lstlisting}[language=Python, caption={FP8 inference with vLLM}]
from vllm import LLM

# Enable FP8 quantization (requires H100/Blackwell or newer)
llm = LLM(
    model="meta-llama/Llama-3-70b-instruct",
    quantization="fp8",           # Use FP8 for both weights and activations
    gpu_memory_utilization=0.90,
    tensor_parallel_size=1,       # 70B FP8 fits on single H100!
)
\end{lstlisting}

% ─────────────────────────────────────────────────────────────
\section{Quantization Accuracy Tradeoffs}

\begin{center}
\begin{tabular}{lrrrr}
  \toprule
  \textbf{Method} & \textbf{Bits} & \textbf{Memory} & \textbf{Speed} &
  \textbf{Accuracy Loss} \\
  \midrule
  BF16 (baseline) & 16 & 140 GB & 1.0$\times$ & 0\% \\
  FP8 & 8 & 70 GB & 1.9$\times$ & $<$0.1\% \\
  INT8 (SmoothQuant) & 8 & 70 GB & 1.7$\times$ & $<$0.5\% \\
  INT4 (AWQ) & 4 & 35 GB & 2.5$\times$ & 1--3\% \\
  INT4 (GPTQ) & 4 & 35 GB & 2.3$\times$ & 1--5\% \\
  FP4 (Blackwell) & 4 & 35 GB & 4.0$\times$ & $<$1\% \\
  \bottomrule
\end{tabular}
\end{center}

\noindent All numbers are approximate for Llama-3 70B on H100.  Accuracy loss
measured as perplexity increase on common benchmarks.

\tipbox{Which Quantization to Choose?}{
\begin{itemize}
  \item \textbf{H100/Blackwell, maximum speed:} FP8 (native hardware support,
        minimal accuracy loss)
  \item \textbf{Memory constrained, quality matters:} AWQ INT4 (best accuracy
        at 4-bit)
  \item \textbf{Maximum compression:} GPTQ INT4 with group size 32
  \item \textbf{Edge/CPU deployment:} GGUF with llama.cpp (Q4\_K\_M is excellent)
\end{itemize}
}

\begin{chapsummary}
Quantization reduces model weight precision to save memory and increase speed.
Lower precision (FP8, INT4) means smaller models, higher throughput, but some
accuracy loss.  GPTQ and AWQ are leading 4-bit post-training quantization methods.
FP8 is now natively supported on H100 and Blackwell GPUs with negligible accuracy
loss.  On Blackwell, FP4 delivers 4$\times$ the throughput of H100 FP8.  Always
benchmark accuracy on your specific task before deploying a quantized model.
\end{chapsummary}

% ============================================================
\chapter{Speculative Decoding}
\label{ch:speculative}
% ============================================================

\begin{learningobjectives}
  \item Understand why speculative decoding speeds up inference
  \item Learn the draft-verify paradigm
  \item Understand acceptance rates and when speculation helps
  \item Apply speculative decoding with vLLM and TensorRT-LLM
\end{learningobjectives}

% ─────────────────────────────────────────────────────────────
\section{The Bottleneck: Sequential Decode Steps}

Recall that decode generates one token at a time.  Each step requires a full
forward pass through the model.  Even if you have a very fast GPU, you cannot
parallelize these steps: token $n+1$ cannot be generated until token $n$ is known.

\analogybox{The Draft and Reviewer}{
Imagine a document that must be approved by an expert (the large model).
The expert is slow but always right.

You hire a fast intern (the small model) to write a draft.  The intern quickly
suggests the next 5 words.  Then the expert reviews all 5 words in a single pass
--- accepting most of them, and correcting any mistakes.

Result: the expert does one review instead of five.  If the intern is correct
4 out of 5 times, the overall speed roughly doubles.

This is \textit{speculative decoding}.
}

% ─────────────────────────────────────────────────────────────
\section{How Speculative Decoding Works}

\index{speculative decoding}
Speculative decoding\index{speculative decoding} uses two models:

\begin{itemize}
  \item \textbf{Draft model:} A small, fast model that quickly generates $k$
        candidate tokens.
  \item \textbf{Target model:} The large, accurate model that verifies all $k$
        tokens in \textbf{one parallel forward pass}.
\end{itemize}

The key insight is that the target model can verify $k$ tokens in parallel
(similar to prefill) rather than generating them sequentially.  If the draft is
mostly correct, we get multiple output tokens per target model forward pass.

\begin{center}
\begin{tikzpicture}[scale=0.9, node distance=0.4cm]
  \tikzstyle{sbox}=[rectangle, draw=AccentOrange, fill=AccentOrange!15,
                    rounded corners=3pt, minimum height=0.7cm,
                    font=\small, inner sep=4pt]
  \tikzstyle{lbox}=[rectangle, draw=NavyBlue, fill=NavyBlue!15,
                    rounded corners=3pt, minimum height=0.7cm,
                    font=\small, inner sep=4pt]

  % Step 1: Draft
  \node[sbox] (d1) {Draft model};
  \node[sbox, right=of d1] (t1) {``The''};
  \node[sbox, right=of t1] (t2) {``sky''};
  \node[sbox, right=of t2] (t3) {``is''};
  \node[sbox, right=of t3] (t4) {``blue''};
  \node[sbox, right=of t4] (t5) {``today''};
  \node[font=\small\bfseries, above=0.1cm of d1] {Step 1: Draft 5 tokens};

  \foreach \a/\b in {d1/t1, t1/t2, t2/t3, t3/t4, t4/t5}
    \draw[-{Stealth}] (\a) -- (\b);

  % Step 2: Verify
  \node[lbox, below=1.2cm of d1] (lm) {Target model (parallel verify)};
  \draw[-{Stealth}] (lm.east) -- ++(5,0)
    node[midway, above, font=\small] {verify all 5 tokens at once};

  % Tick/cross
  \foreach \x/\res/\col in {
    3.2/\checkmark/AccentGreen, 4.6/\checkmark/AccentGreen,
    6.0/\checkmark/AccentGreen, 7.4/\checkmark/AccentGreen,
    8.8/$\times$/red}
  {
    \node[font=\large\bfseries, text=\col] at (\x, -1.8) {\res};
  }
  \node[font=\small\bfseries, below=3cm of d1]
    {4 tokens accepted, 1 rejected (acceptance rate = 80\%)};
\end{tikzpicture}
\end{center}

\subsection{N-gram and EAGLE Speculation}

Beyond a separate small model, modern systems use:

\begin{itemize}
  \item \textbf{N-gram speculation:} Predict next tokens based on n-gram statistics
        from the current context.  Fast and requires no extra model.
  \item \textbf{EAGLE/EAGLE-3:} Uses a shallow draft head that shares the target
        model's hidden states.  Achieves 3$\times$ speedup on typical text.
  \item \textbf{Medusa:} Multiple parallel draft heads on the same model backbone.
\end{itemize}

\begin{lstlisting}[language=Python, caption={Speculative decoding with vLLM}]
from vllm import LLM, SamplingParams

# Method 1: Use a small draft model
llm = LLM(
    model="meta-llama/Llama-3-70b-instruct",        # Target model
    speculative_model="meta-llama/Llama-3-8b-instruct",  # Draft model
    num_speculative_tokens=5,   # Generate 5 draft tokens per step
    tensor_parallel_size=2,
)

# Method 2: N-gram based speculation (no extra model needed)
llm_ngram = LLM(
    model="meta-llama/Llama-3-70b-instruct",
    speculative_model="[ngram]",     # Built-in n-gram drafter
    ngram_prompt_lookup_max=4,       # Look back 4 tokens for patterns
    num_speculative_tokens=5,
    tensor_parallel_size=2,
)

result = llm.generate(
    ["Write a Python function to sort a list:"],
    SamplingParams(temperature=0.0, max_tokens=200)  # Greedy is ideal for speculation
)
\end{lstlisting}

\notebox{When Speculation Helps Most}{
Speculative decoding helps most when:
\begin{itemize}
  \item Output is predictable (code, structured text, repetitive responses)
  \item Batch size is small (low utilization otherwise)
  \item Latency matters more than throughput
\end{itemize}
Speculative decoding \textit{hurts} throughput at high batch sizes because it adds
verification overhead.  At a batch size of 32+, continuous batching is usually
better than speculation.
}

\begin{chapsummary}
Speculative decoding uses a fast draft model to propose multiple tokens, which the
target model verifies in parallel.  When the draft is mostly correct, this delivers
2--4$\times$ lower latency for single-request serving.  EAGLE-3 and n-gram
speculation are the preferred methods in 2025.  Speculation is most beneficial for
low-batch, latency-sensitive use cases, not high-throughput batch serving.
\end{chapsummary}

% ============================================================
\chapter{Inference Runtimes: vLLM, TensorRT-LLM, and SGLang}
\label{ch:runtimes}
% ============================================================

\begin{learningobjectives}
  \item Understand the architecture of the three leading inference runtimes
  \item Compare vLLM, TensorRT-LLM, and SGLang across key dimensions
  \item Know which runtime to choose for a given use case
  \item Deploy a model with vLLM and a REST API
\end{learningobjectives}

% ─────────────────────────────────────────────────────────────
\section{The Inference Runtime Stack}

An \keyterm{inference runtime}\index{inference runtime} is the software layer that
takes a trained model and serves it efficiently to users.  Think of it as the
kitchen of a restaurant: the model is the recipe, and the runtime is the kitchen
with all its equipment, scheduling, and organization.

Every production runtime must handle:
\begin{itemize}
  \item \textbf{Model loading:} Reading weights from disk, placing on GPUs
  \item \textbf{Request queuing:} Accepting HTTP/gRPC requests from users
  \item \textbf{Batching:} Grouping multiple requests together for efficiency
  \item \textbf{Scheduling:} Deciding which requests to process next
  \item \textbf{Memory management:} Allocating/freeing KV cache for each request
  \item \textbf{Output:} Streaming tokens back to users as they are generated
\end{itemize}

% ─────────────────────────────────────────────────────────────
\section{vLLM: The Production Standard}
\index{vLLM}

\keyterm{vLLM} (UC Berkeley, 2023, now a PyTorch Foundation project) is the
most widely adopted open-source inference engine.  Its core innovations:

\begin{itemize}
  \item \textbf{PagedAttention:} Eliminates KV cache fragmentation (Chapter~\ref{ch:kvcache})
  \item \textbf{Continuous batching:} New requests join as soon as a slot frees up
  \item \textbf{Prefix caching:} Reuses KV state for shared prefixes
  \item \textbf{OpenAI-compatible API:} Drop-in replacement for OpenAI API
\end{itemize}

Typical performance: 120--160 requests/second, 50--80ms TTFT with 100 concurrent
users for Llama-class models on H100.

\subsection{Continuous Batching}
\index{continuous batching}

Traditional batching waits for a batch to be full before starting.  This is like
a bus that only departs when every seat is full --- efficient for the bus, but
terrible for passengers who arrived early.

\keyterm{Continuous batching} (also called iteration-level scheduling) allows
new requests to join the batch at every decode step.  When request A finishes at
step 50, a new request immediately fills that slot.

\begin{center}
\begin{tikzpicture}[scale=0.75]
  % Static batching (bad)
  \node[font=\small\bfseries\color{NavyBlue}] at (3,6.5) {Static Batching};
  \draw[fill=NavyBlue!10, draw=NavyBlue, rounded corners=3pt]
    (0,5.5) rectangle (6,6.3);
  \node[font=\tiny] at (3,5.9) {Req A (20 steps), Req B (20 steps), Req C (20 steps)};
  \draw[fill=AccentOrange!20, draw=AccentOrange, rounded corners=3pt]
    (0,4.5) rectangle (6,5.3);
  \node[font=\tiny\color{AccentOrange}] at (3,4.9) {WAIT for next batch...};
  \draw[fill=NavyBlue!10, draw=NavyBlue, rounded corners=3pt]
    (0,3.5) rectangle (6,4.3);
  \node[font=\tiny] at (3,3.9) {Req D, E, F (waiting all this time!)};

  % Arrow
  \draw[-{Stealth}] (7,5) -- (8,5);

  % Continuous batching (good)
  \node[font=\small\bfseries\color{AccentGreen}] at (12,6.5)
    {Continuous Batching};
  \foreach \y/\col/\lbl in {
    5.8/NavyBlue/Req A (finishes early),
    4.9/NavyBlue/Req B,
    4.0/NavyBlue/Req C,
    3.1/AccentGreen/Req D joins immediately!}
  {
    \draw[fill=\col!20, draw=\col, rounded corners=2pt]
      (8.5,\y-0.35) rectangle (15.5,\y+0.35);
    \node[font=\tiny] at (12,\y) {\lbl};
  }
\end{tikzpicture}
\end{center}

\subsection{Running vLLM}

\begin{lstlisting}[style=bashstyle, caption={Starting the vLLM OpenAI-compatible server}]
# Install vLLM
pip install vllm

# Start the server (OpenAI-compatible API on port 8000)
python -m vllm.entrypoints.openai.api_server \
    --model meta-llama/Llama-3-70b-instruct \
    --tensor-parallel-size 2 \          # Use 2 GPUs
    --gpu-memory-utilization 0.90 \     # Use 90% of GPU memory for KV cache
    --enable-prefix-caching \           # Enable prefix caching
    --kv-cache-dtype fp8 \              # FP8 KV cache
    --max-num-seqs 256 \                # Max concurrent sequences
    --port 8000
\end{lstlisting}

\begin{lstlisting}[language=Python, caption={Querying the vLLM server}]
import openai

# vLLM is OpenAI API compatible
client = openai.OpenAI(
    base_url="http://localhost:8000/v1",
    api_key="not-needed"                # vLLM doesn't require a key
)

response = client.chat.completions.create(
    model="meta-llama/Llama-3-70b-instruct",
    messages=[
        {"role": "system", "content": "You are a helpful assistant."},
        {"role": "user", "content": "Explain KV cache in 2 sentences."}
    ],
    max_tokens=200,
    stream=True                         # Enable streaming
)

for chunk in response:
    if chunk.choices[0].delta.content:
        print(chunk.choices[0].delta.content, end="", flush=True)
\end{lstlisting}

% ─────────────────────────────────────────────────────────────
\section{TensorRT-LLM: Maximum NVIDIA Performance}
\index{TensorRT-LLM}

\keyterm{TensorRT-LLM} (NVIDIA) compiles a model into a highly optimized engine
specifically for a particular GPU, precision, and sequence configuration.

Key advantages:
\begin{itemize}
  \item Lowest latency at low concurrency (35--50ms TTFT vs.\ vLLM's 50--80ms)
  \item Deep GPU kernel fusion (multiple operations merged into one kernel)
  \item FP8/FP4 native support for Blackwell
  \item Disaggregated serving via NVIDIA Dynamo
\end{itemize}

Key disadvantages:
\begin{itemize}
  \item Complex engine compilation (can take hours per model)
  \item NVIDIA-only (no AMD, Intel)
  \item Requires re-compilation when model or GPU changes
\end{itemize}

\begin{lstlisting}[style=bashstyle, caption={Building and running a TensorRT-LLM engine}]
# Install TensorRT-LLM
pip install tensorrt_llm

# Step 1: Convert model to TensorRT-LLM checkpoint format
python convert_checkpoint.py \
    --model_dir ./llama-3-70b \
    --output_dir ./trt-checkpoint \
    --dtype float16

# Step 2: Build the TRT engine (this compiles for your specific GPU)
trtllm-build \
    --checkpoint_dir ./trt-checkpoint \
    --output_dir ./trt-engine \
    --max_input_len 2048 \
    --max_output_len 512 \
    --max_batch_size 32 \
    --tp_size 2                  # Tensor parallel across 2 GPUs

# Step 3: Run inference
python run.py \
    --engine_dir ./trt-engine \
    --max_output_len 100 \
    --input_text "Explain GPU memory hierarchy:"
\end{lstlisting}

% ─────────────────────────────────────────────────────────────
\section{SGLang: Structured Generation Language}
\index{SGLang}

\keyterm{SGLang} (Zheng et al., 2024) introduces \textit{RadixAttention} --- an
automatic KV cache reuse system that builds a radix tree of all cached prefixes
and shares them optimally across requests.

SGLang excels for:
\begin{itemize}
  \item Long-document processing (PDF analysis, RAG)
  \item Conversation-style traffic with long shared history
  \item Multi-turn agents that reuse the same system context
  \item Structured generation (JSON output, grammars)
\end{itemize}

\begin{lstlisting}[style=bashstyle, caption={Starting SGLang server}]
pip install sglang

# Launch SGLang server
python -m sglang.launch_server \
    --model-path meta-llama/Llama-3-70b-instruct \
    --tp 2 \                     # Tensor parallel size
    --mem-fraction-static 0.82 \
    --enable-flashinfer          # Use FlashInfer attention backend
\end{lstlisting}

% ─────────────────────────────────────────────────────────────
\section{Runtime Comparison: The Decision Guide}

\begin{longtable}{@{}p{3cm}p{3cm}p{3cm}p{3cm}@{}}
  \toprule
  \textbf{Criterion} & \textbf{vLLM} & \textbf{TRT-LLM} & \textbf{SGLang} \\
  \midrule
  \endfirsthead
  \toprule
  \textbf{Criterion} & \textbf{vLLM} & \textbf{TRT-LLM} & \textbf{SGLang} \\
  \midrule
  \endhead
  Setup time & Minutes & Hours/days & Minutes \\
  TTFT (low concurrency) & 50--80ms & 35--50ms & 50--80ms \\
  Throughput (high concurrency) & Excellent & Excellent & Excellent \\
  Hardware & Any (NVIDIA/AMD) & NVIDIA only & NVIDIA/AMD \\
  Prefix caching & Good & Good & Excellent (RadixAttention) \\
  Model support & Widest & NVIDIA-curated & Growing \\
  API compatibility & OpenAI & Custom/Triton & OpenAI \\
  Best for & General production & Max NVIDIA perf & Long context/RAG \\
  \bottomrule
  \caption{Inference Runtime Comparison (2025)}
\end{longtable}

\begin{chapsummary}
vLLM is the go-to for most production deployments: easy setup, OpenAI-compatible,
excellent throughput.  TensorRT-LLM achieves the lowest latency on NVIDIA hardware
but requires compilation effort.  SGLang leads for long-context workloads with
RadixAttention.  All three support modern features: FP8, continuous batching,
prefix caching.  Choose based on your traffic shape, ops maturity, and latency SLAs.
\end{chapsummary}

% ============================================================
\chapter{Distributed Inference: Parallelism Strategies}
\label{ch:distributed}
% ============================================================

\begin{learningobjectives}
  \item Understand why a single GPU is often not enough
  \item Learn tensor parallelism, pipeline parallelism, and sequence parallelism
  \item Understand NVLink and its role in multi-GPU scaling
  \item Apply tensor parallelism with vLLM
\end{learningobjectives}

% ─────────────────────────────────────────────────────────────
\section{When One GPU Is Not Enough}

A Llama-3 70B model in FP16 requires 140 GB of memory --- far more than a single
H100's 80 GB.  Even smaller models may benefit from multiple GPUs for throughput.

Distributing a model across multiple GPUs requires a \textbf{parallelism strategy}.
There are three main approaches:

% ─────────────────────────────────────────────────────────────
\section{Tensor Parallelism}
\index{tensor parallelism}

\keyterm{Tensor parallelism} (TP) splits individual matrices \textit{across} GPUs.

\analogybox{The Heavy Book}{
Imagine a very large book that is too heavy for one person to carry.
Instead of one person carrying the whole book, two people each carry half ---
one carries pages 1--500, the other carries pages 501--1000.

Tensor parallelism works similarly.  A large weight matrix (say, $8192 \times 8192$)
is split across 2 GPUs: each GPU stores a $4096 \times 8192$ slice.

Both GPUs compute in parallel, then combine their partial results.  The
communication (combining results) uses NVLink for very high bandwidth.
}

With TP-4 (4 GPUs), each GPU holds 1/4 of each weight matrix.  Each forward
pass requires an \textbf{AllReduce} communication operation to combine results
across all 4 GPUs.

\begin{center}
\begin{tikzpicture}[scale=0.85]
  % Weight matrix
  \node[font=\small\bfseries] at (-2,2) {Weight $W$};
  \draw[fill=NavyBlue!20, draw=NavyBlue] (-3.5,0) rectangle (-0.5,4);

  % Split arrow
  \draw[-{Stealth}, line width=1.5pt, AccentOrange] (-0.5,2) -- (0.5,2)
    node[midway,above,font=\tiny] {split};

  % GPU 0
  \draw[fill=AccentGreen!20, draw=AccentGreen, rounded corners=3pt]
    (0.7, 2.0) rectangle (3.7, 4.0);
  \node[font=\small\bfseries\color{AccentGreen}] at (2.2,3.0) {GPU 0: $W[:n/2, :]$};

  % GPU 1
  \draw[fill=NavyBlue!20, draw=NavyBlue, rounded corners=3pt]
    (0.7, -0.2) rectangle (3.7, 1.8);
  \node[font=\small\bfseries\color{NavyBlue}] at (2.2,0.8) {GPU 1: $W[n/2:, :]$};

  % AllReduce
  \draw[-{Stealth}, AccentOrange, line width=1.5pt] (3.7,3.0) -- (5.0,2.0)
    node[midway,above right,font=\tiny] {AllReduce};
  \draw[-{Stealth}, AccentOrange, line width=1.5pt] (3.7,0.8) -- (5.0,2.0);
  \draw[fill=AccentOrange!20, draw=AccentOrange, rounded corners=3pt]
    (5.0,1.5) rectangle (7.5,2.5);
  \node[font=\small] at (6.25,2.0) {Combined};
\end{tikzpicture}
\end{center}

\subsection{Launching Tensor-Parallel vLLM}

\begin{lstlisting}[style=bashstyle, caption={Running vLLM with tensor parallelism}]
# 70B model across 2 H100s with TP=2
python -m vllm.entrypoints.openai.api_server \
    --model meta-llama/Llama-3-70b-instruct \
    --tensor-parallel-size 2        # Split across 2 GPUs

# 405B model across 8 H100s with TP=8
python -m vllm.entrypoints.openai.api_server \
    --model meta-llama/Llama-3-405b-instruct \
    --tensor-parallel-size 8        # Split across 8 GPUs
\end{lstlisting}

% ─────────────────────────────────────────────────────────────
\section{Pipeline Parallelism}
\index{pipeline parallelism}

\keyterm{Pipeline parallelism} (PP) assigns different model \textit{layers} to
different GPUs.  Layers 0--19 run on GPU 0, layers 20--39 run on GPU 1, etc.

\analogybox{The Assembly Line}{
In a car factory, one station installs the engine, another the doors, another the
paint.  Each station works on a different car simultaneously.

In pipeline parallelism, GPU 0 finishes processing layers 0--19 for batch 1 and
passes the intermediate activations to GPU 1.  While GPU 1 works on batch 1,
GPU 0 starts processing batch 2.

The pipeline keeps all GPUs busy, like a factory assembly line.
}

Pipeline parallelism communicates via a simpler point-to-point activation transfer
(not AllReduce), making it more efficient across slower interconnects
(e.g., InfiniBand, PCIe) than tensor parallelism.

% ─────────────────────────────────────────────────────────────
\section{NVLink: The Fast Lane Between GPUs}
\index{NVLink}

Communication between GPUs is the bottleneck in distributed inference.  NVLink is
NVIDIA's high-speed GPU interconnect, far faster than PCIe:

\begin{center}
\begin{tabular}{lll}
  \toprule
  \textbf{Interconnect} & \textbf{Bandwidth (each dir.)} & \textbf{Use Case} \\
  \midrule
  PCIe 5.0 $\times$16 & 64 GB/s & CPU--GPU \\
  NVLink 4 (H100) & 900 GB/s & GPU--GPU \\
  NVLink 5 (B200) & 1800 GB/s & GPU--GPU \\
  NVSwitch (GB200 NVL72) & 57.6 TB/s total & Rack-scale \\
  \bottomrule
\end{tabular}
\end{center}

With NVLink 4 on a DGX H100 system (8 GPUs), the AllReduce bandwidth for
tensor parallelism is \textbf{28$\times$} faster than a PCIe-only configuration.
This is why TP works well intra-node but poorly across nodes.

\tipbox{The Parallelism Recipe}{
\begin{itemize}
  \item \textbf{TP within a node} (GPUs on same server, NVLink connected)
  \item \textbf{PP across nodes} (different servers, InfiniBand connected)
  \item \textbf{TP=8, PP=N/8 for very large models} (405B+)
\end{itemize}
}

% ─────────────────────────────────────────────────────────────
\section{Disaggregated Prefill and Decode}
\index{disaggregated serving}

An emerging production pattern separates the prefill and decode phases onto
\textbf{different hardware}:

\begin{itemize}
  \item \textbf{Prefill GPUs:} Optimized for high compute (processing long prompts)
  \item \textbf{Decode GPUs:} Optimized for memory bandwidth (fast token generation)
\end{itemize}

NVIDIA Dynamo (announced GTC 2025) is an orchestration layer that implements this
pattern on top of vLLM, TensorRT-LLM, or SGLang.  By routing prefill and decode to
specialized hardware, it can significantly reduce queue depth and improve both TTFT
and throughput simultaneously.

\begin{chapsummary}
Large models require multi-GPU deployment.  Tensor parallelism splits weight
matrices across GPUs --- ideal for intra-node with NVLink.  Pipeline parallelism
assigns layers to GPUs --- better for inter-node.  NVLink provides the
high-bandwidth interconnect that makes tensor parallelism practical.  Disaggregated
prefill/decode is an emerging frontier that dedicates different hardware to the
two phases of inference.
\end{chapsummary}

% ============================================================
\chapter{Observability: Monitoring Inference Systems}
\label{ch:observability}
% ============================================================

\begin{learningobjectives}
  \item Understand the key metrics to monitor in a production inference system
  \item Learn to use NVIDIA profiling tools: nvtop, nsys, ncu
  \item Set up Prometheus and Grafana for inference monitoring
  \item Diagnose common performance bottlenecks from metrics
\end{learningobjectives}

% ─────────────────────────────────────────────────────────────
\section{The Three Pillars of Observability}

A production inference system must be observable:

\begin{itemize}
  \item \textbf{Metrics:} Numerical measurements over time (latency p99, GPU util\%)
  \item \textbf{Logs:} Events with timestamps (errors, request traces)
  \item \textbf{Traces:} Request-level timing breakdown (prefill time, queue time)
\end{itemize}

% ─────────────────────────────────────────────────────────────
\section{Key Metrics to Monitor}

\begin{longtable}{@{}p{3.5cm}p{4cm}p{5cm}@{}}
  \toprule
  \textbf{Metric} & \textbf{What It Measures} & \textbf{Action if Bad} \\
  \midrule
  \endfirsthead
  \toprule
  \textbf{Metric} & \textbf{What It Measures} & \textbf{Action if Bad} \\
  \midrule
  \endhead
  GPU utilization & \% time GPU is active & $<$60\%: add batching; increase load \\
  GPU memory \% & VRAM in use & $>$95\%: reduce batch size or add GPUs \\
  TTFT (p50/p95/p99) & Latency to first token & High: check queue depth, prefill load \\
  TPOT (p50/p95) & Latency per output token & High: check decode batch size \\
  Throughput (tok/s) & Tokens generated per second & Low: check GPU util, batching \\
  Request queue depth & Requests waiting & $>$0 for $>$1s: scale up \\
  KV cache usage \% & Cache pool utilization & Near 100\%: risk of OOM eviction \\
  Token generation rate & New tokens/s overall & Declining: memory pressure \\
  \bottomrule
  \caption{Production Inference Monitoring Metrics}
\end{longtable}

% ─────────────────────────────────────────────────────────────
\section{GPU Profiling Tools}
\index{GPU profiling}

\subsection{nvtop: Real-Time GPU Monitor}

\begin{lstlisting}[style=bashstyle, caption={Using nvtop for GPU monitoring}]
# Install nvtop
apt-get install nvtop

# Run nvtop
nvtop
# Shows: GPU util%, memory, power, temperature for all GPUs
# Press F2 to configure display, q to quit
\end{lstlisting}

\subsection{Nsight Systems: Timeline Profiling}
\index{Nsight Systems}

\begin{lstlisting}[style=bashstyle, caption={Profiling with Nsight Systems}]
# Profile a vLLM inference run
nsys profile \
    --trace cuda,nvtx,osrt \
    --output inference_profile \
    python inference_benchmark.py

# Open in Nsight Systems GUI
nsys-ui inference_profile.nsys-rep
# Shows: kernel timelines, memory transfers, CPU/GPU overlap
\end{lstlisting}

\subsection{Nsight Compute: Kernel-Level Analysis}
\index{Nsight Compute}

\begin{lstlisting}[style=bashstyle, caption={Per-kernel profiling with ncu}]
# Profile specific CUDA kernels (use on small workloads)
ncu --set full \
    --kernel-name "flash_fwd_kernel" \
    --export kernel_profile \
    python run_attention.py

# Key metrics to look for:
# - SM active cycles: is GPU compute actually busy?
# - Memory throughput: are we hitting bandwidth limits?
# - Warp efficiency: are warps diverging?
\end{lstlisting}

% ─────────────────────────────────────────────────────────────
\section{Prometheus and Grafana for Inference}
\index{Prometheus}\index{Grafana}

vLLM exposes metrics natively at \texttt{/metrics} in Prometheus format:

\begin{lstlisting}[style=bashstyle, caption={Starting vLLM with metrics}]
python -m vllm.entrypoints.openai.api_server \
    --model meta-llama/Llama-3-70b-instruct \
    --tensor-parallel-size 2 \
    --port 8000 \
    --metrics-enabled          # Enable Prometheus metrics
# Metrics available at http://localhost:8000/metrics
\end{lstlisting}

\begin{lstlisting}[style=yamlstyle, caption={Prometheus scrape config for vLLM}]
# prometheus.yml
scrape_configs:
  - job_name: 'vllm'
    static_configs:
      - targets: ['localhost:8000']
    metrics_path: '/metrics'
    scrape_interval: 5s

# Key vLLM metrics:
# vllm:num_requests_running        - active requests
# vllm:gpu_cache_usage_perc        - KV cache % used
# vllm:e2e_request_latency_seconds - end-to-end latency histogram
# vllm:generation_tokens_total     - tokens generated
\end{lstlisting}

\begin{chapsummary}
Production inference requires continuous monitoring.  Key metrics are GPU
utilization, KV cache usage, TTFT, TPOT, and throughput.  NVIDIA tools (nvtop,
nsys, ncu) provide GPU-level insight.  vLLM's Prometheus metrics enable Grafana
dashboards for operational monitoring.  High queue depth means scale up.  Low GPU
utilization means increase batch size or adjust scheduling.
\end{chapsummary}

% ============================================================
\chapter{Production Deployment with Kubernetes}
\label{ch:kubernetes}
% ============================================================

\begin{learningobjectives}
  \item Understand how to containerize and deploy inference servers
  \item Configure GPU resources in Kubernetes
  \item Implement auto-scaling based on latency metrics
  \item Set up health checks and rolling updates for zero-downtime deployments
\end{learningobjectives}

% ─────────────────────────────────────────────────────────────
\section{Containerizing a vLLM Server}

\begin{lstlisting}[style=bashstyle, caption={Dockerfile for a vLLM inference server}]
# Dockerfile
FROM nvcr.io/nvidia/cuda:12.4.0-devel-ubuntu22.04

# Install Python and dependencies
RUN apt-get update && apt-get install -y python3-pip python3-dev && \
    pip install vllm==0.7.0 huggingface_hub

# Pre-download the model (or mount as a volume in production)
RUN python3 -c "from huggingface_hub import snapshot_download; \
    snapshot_download('meta-llama/Llama-3-8b-instruct')"

# Expose the API port
EXPOSE 8000

# Start the vLLM server
CMD ["python3", "-m", "vllm.entrypoints.openai.api_server", \
     "--model", "meta-llama/Llama-3-8b-instruct", \
     "--port", "8000", \
     "--gpu-memory-utilization", "0.90"]
\end{lstlisting}

% ─────────────────────────────────────────────────────────────
\section{Kubernetes Deployment}

\begin{lstlisting}[style=yamlstyle, caption={Kubernetes Deployment for vLLM}]
# vllm-deployment.yaml
apiVersion: apps/v1
kind: Deployment
metadata:
  name: vllm-llama3-8b
  labels:
    app: vllm-inference
spec:
  replicas: 2                  # Start with 2 replicas
  selector:
    matchLabels:
      app: vllm-inference
  template:
    metadata:
      labels:
        app: vllm-inference
    spec:
      containers:
      - name: vllm
        image: my-registry/vllm-llama3:latest
        ports:
        - containerPort: 8000
        resources:
          limits:
            nvidia.com/gpu: "1"      # 1 GPU per pod
            memory: "40Gi"
            cpu: "8"
          requests:
            nvidia.com/gpu: "1"
            memory: "32Gi"
            cpu: "4"
        env:
        - name: CUDA_VISIBLE_DEVICES
          value: "0"
        livenessProbe:
          httpGet:
            path: /health
            port: 8000
          initialDelaySeconds: 60
          periodSeconds: 30
        readinessProbe:
          httpGet:
            path: /health
            port: 8000
          initialDelaySeconds: 90
          periodSeconds: 10
---
apiVersion: v1
kind: Service
metadata:
  name: vllm-service
spec:
  selector:
    app: vllm-inference
  ports:
  - port: 80
    targetPort: 8000
  type: LoadBalancer
\end{lstlisting}

% ─────────────────────────────────────────────────────────────
\section{Horizontal Pod Autoscaler}

\begin{lstlisting}[style=yamlstyle, caption={HPA based on custom latency metric}]
# hpa-vllm.yaml
apiVersion: autoscaling/v2
kind: HorizontalPodAutoscaler
metadata:
  name: vllm-hpa
spec:
  scaleTargetRef:
    apiVersion: apps/v1
    kind: Deployment
    name: vllm-llama3-8b
  minReplicas: 2
  maxReplicas: 10
  metrics:
  - type: Pods
    pods:
      metric:
        name: vllm_queue_depth    # Scale when queue is growing
      target:
        type: AverageValue
        averageValue: "5"         # Scale up when avg queue > 5 requests
\end{lstlisting}

\begin{lstlisting}[style=bashstyle, caption={Deploying to Kubernetes}]
# Apply deployment
kubectl apply -f vllm-deployment.yaml

# Check pod status
kubectl get pods -l app=vllm-inference

# Check GPU allocation
kubectl describe pod vllm-llama3-8b-<pod-id>

# View logs
kubectl logs -f deployment/vllm-llama3-8b

# Check service
kubectl get service vllm-service
\end{lstlisting}

\begin{chapsummary}
Production inference runs on Kubernetes with GPU-aware pods.  Containerize with
NVIDIA CUDA base images and expose a health endpoint.  Specify
\texttt{nvidia.com/gpu} resource limits.  Use HPA with custom metrics (queue
depth, latency) for autoscaling.  Set readiness probes with long initial delays
--- model loading takes time.  Use rolling updates for zero-downtime deployments.
\end{chapsummary}

% ============================================================
%  PART II: ADVANCED TOPICS
% ============================================================
\part{Advanced Topics}

% ============================================================
\chapter{Long Context, Multimodal, and Agentic Inference}
\label{ch:advanced}
% ============================================================

\begin{learningobjectives}
  \item Understand the unique challenges of long-context inference
  \item Learn how multimodal models handle images, audio, and video
  \item Understand agentic inference patterns and their infrastructure needs
\end{learningobjectives}

\section{Long-Context Inference}
\index{long context}

Models with 128K--1M token context windows introduce extreme memory and compute
challenges.  The KV cache for a single 128K-token conversation with Llama-3 70B
reaches:

\[
  2 \times 80 \times 64 \times 128 \times 131072 \times 2 \approx \mathbf{344\text{ GB}}
\]

This dwarfs even the GPU's HBM capacity.  Techniques for long-context efficiency:

\begin{itemize}
  \item \textbf{Sparse attention:} Only attend to a subset of past tokens
        (e.g., StreamingLLM's ``attention sinks'')
  \item \textbf{KV cache quantization to INT2/INT4:} Sub-4-bit KV compression
  \item \textbf{KV cache offloading:} Spill to CPU memory or NVMe SSD
  \item \textbf{Chunked prefill:} Process long prompts in chunks to control
        peak memory
\end{itemize}

\begin{lstlisting}[language=Python, caption={Chunked prefill for long contexts in vLLM}]
from vllm import LLM, SamplingParams

llm = LLM(
    model="meta-llama/Llama-3-70b-instruct",
    max_model_len=131072,          # 128K context window
    enable_chunked_prefill=True,   # Process long prompts in chunks
    max_num_batched_tokens=2048,   # Process 2K tokens per prefill chunk
    tensor_parallel_size=4,        # Need 4 GPUs for 128K KV cache
)
\end{lstlisting}

\section{Multimodal Inference}
\index{multimodal inference}

Models like LLaVA, Qwen-VL, and InternVL accept image inputs alongside text.
The inference pipeline has an extra step: image encoding.

\begin{enumerate}
  \item \textbf{Vision encoder:} A ViT (Vision Transformer) encodes the image
        into feature tokens --- typically 256--1024 tokens per image.
  \item \textbf{Projection:} A linear layer maps image features to the LLM's
        token space.
  \item \textbf{LLM decode:} The LLM generates text given the image tokens
        and text prompt.
\end{enumerate}

The bottleneck is that each image adds many tokens to the effective sequence
length, inflating KV cache requirements significantly.

\section{Agentic Inference}
\index{agentic inference}

AI agents run multiple LLM calls in a loop: plan, act, observe, repeat.
The infrastructure implications:

\begin{itemize}
  \item \textbf{Long conversations:} Context accumulates over many steps.
        Efficient prefix caching is essential.
  \item \textbf{Tool calling:} The model generates structured JSON for tool calls.
        Low latency per step matters for responsiveness.
  \item \textbf{Parallel function calls:} Some frameworks call multiple tools
        simultaneously.  The serving infrastructure must handle bursty traffic.
  \item \textbf{Request tracing:} Agents make many sub-requests.  Tracing is
        essential for debugging and cost attribution.
\end{itemize}

\begin{chapsummary}
Long-context inference strains KV cache memory and requires chunked prefill,
sparse attention, or quantized KV.  Multimodal models add a vision encoder step
that contributes many extra tokens.  Agentic workloads create bursty, highly
correlated requests that benefit from aggressive prefix caching and low per-call
latency.
\end{chapsummary}

% ============================================================
%  APPENDICES
% ============================================================
\appendix

% ============================================================
\chapter{CUDA Cheat Sheet}
\label{app:cuda}
% ============================================================

\section*{Essential CUDA API}

\begin{lstlisting}[style=cstyle, caption={CUDA Cheat Sheet}]
/* ── Memory ─────────────────────────────────────────────── */
cudaMalloc(&ptr, bytes);             // Allocate GPU memory
cudaFree(ptr);                       // Free GPU memory
cudaMemcpy(dst, src, bytes, kind);   // Transfer: kind =
  // cudaMemcpyHostToDevice | cudaMemcpyDeviceToHost | DeviceToDevice
cudaMemset(ptr, value, bytes);       // Fill GPU memory

/* ── Async memory (for overlap with compute) ───────────── */
cudaMallocAsync(&ptr, bytes, stream);
cudaFreeAsync(ptr, stream);
cudaMemcpyAsync(dst, src, bytes, kind, stream);

/* ── Kernel launch ─────────────────────────────────────── */
kernel<<<grid, block, shared_mem, stream>>>(args...);
// grid  = dim3(gridX, gridY, gridZ)
// block = dim3(blockX, blockY, blockZ)
// max threads per block: 1024

/* ── Thread indexing inside kernel ────────────────────── */
int idx = blockIdx.x * blockDim.x + threadIdx.x;
int idy = blockIdx.y * blockDim.y + threadIdx.y;

/* ── Synchronization ───────────────────────────────────── */
cudaDeviceSynchronize();             // Wait for ALL GPU work
cudaStreamSynchronize(stream);       // Wait for one stream
__syncthreads();                     // Sync threads within block (in kernel)

/* ── Error checking (always do this!) ─────────────────── */
#define CUDA_CHECK(call) { \
    cudaError_t err = call; \
    if (err != cudaSuccess) { \
        fprintf(stderr, "CUDA error: %s at %s:%d\n", \
                cudaGetErrorString(err), __FILE__, __LINE__); \
        exit(1); \
    } \
}
CUDA_CHECK(cudaMalloc(&ptr, bytes));

/* ── Streams ─────────────────────────────────────────── */
cudaStream_t stream;
cudaStreamCreate(&stream);
cudaStreamDestroy(stream);

/* ── Device query ─────────────────────────────────────── */
cudaDeviceProp prop;
cudaGetDeviceProperties(&prop, device_id);
// prop.name, prop.totalGlobalMem, prop.multiProcessorCount
// prop.maxThreadsPerBlock, prop.sharedMemPerBlock
\end{lstlisting}

\section*{Performance Rules of Thumb}

\begin{itemize}
  \item Use at least 128 threads per block; 256 is usually optimal
  \item Ensure memory accesses are coalesced (stride-1 access by consecutive threads)
  \item Avoid warp divergence in performance-critical paths
  \item Use shared memory for data reused $>$2 times within a block
  \item Overlap compute and memory transfers using CUDA streams
  \item Profile with nsys before optimizing --- never guess the bottleneck
\end{itemize}

% ============================================================
\chapter{Linux Commands for AI Engineers}
\label{app:linux}
% ============================================================

\begin{lstlisting}[style=bashstyle, caption={Essential Linux commands for AI engineers}]
# ── GPU monitoring ──────────────────────────────────────────
nvidia-smi                          # GPU status (util, mem, temp)
nvidia-smi dmon -s u                # Live utilization every 1s
nvidia-smi --query-gpu=name,memory.total,memory.free \
    --format=csv                    # GPU memory info
watch -n 1 nvidia-smi              # Auto-refresh every 1s
nvtop                               # Interactive GPU monitor

# ── Process management ──────────────────────────────────────
htop                                # Interactive process viewer
ps aux | grep python               # Find Python processes
kill -9 <pid>                       # Force kill process
nohup python server.py &           # Run in background

# ── Disk and memory ─────────────────────────────────────────
df -h                               # Disk space
du -sh /models/*                   # Model directory sizes
free -h                             # RAM usage
cat /proc/meminfo | grep HugePages # Huge pages (for pinned memory)

# ── Network ─────────────────────────────────────────────────
ss -tlnp | grep 8000               # Check if port 8000 is open
curl http://localhost:8000/health   # Test health endpoint
curl http://localhost:8000/v1/models # List loaded models

# ── Benchmarking ────────────────────────────────────────────
# Time a command
time python inference_benchmark.py

# Test vLLM throughput
python -m vllm.entrypoints.benchmark_throughput \
    --model meta-llama/Llama-3-8b-instruct \
    --input-len 512 --output-len 128 \
    --num-prompts 1000

# ── Environment ─────────────────────────────────────────────
echo $CUDA_VISIBLE_DEVICES         # Which GPUs are visible?
export CUDA_VISIBLE_DEVICES=0,1    # Restrict to GPU 0 and 1
env | grep CUDA                    # All CUDA env variables

# ── Docker/Container ────────────────────────────────────────
docker run --gpus all \            # Use all GPUs in container
    -p 8000:8000 \
    my-vllm-image

# ── Tmux for long-running jobs ──────────────────────────────
tmux new -s inference              # New session named "inference"
tmux attach -t inference           # Reattach to session
# Ctrl+b, d to detach (session keeps running)
\end{lstlisting}

% ============================================================
\chapter{Benchmarking Toolkit}
\label{app:benchmark}
% ============================================================

\begin{lstlisting}[language=Python, caption={Comprehensive inference benchmarking script}]
"""
AI Inference Benchmarking Toolkit
Measures: TTFT, TPOT, throughput, memory usage
"""
import time
import asyncio
import aiohttp
import statistics
import json
from dataclasses import dataclass, field
from typing import List

@dataclass
class BenchmarkResult:
    ttft_ms: List[float] = field(default_factory=list)
    tpot_ms: List[float] = field(default_factory=list)
    total_tokens: List[int] = field(default_factory=list)
    errors: int = 0

    def summary(self):
        return {
            "TTFT_p50_ms":  statistics.median(self.ttft_ms),
            "TTFT_p95_ms":  statistics.quantiles(self.ttft_ms, n=20)[18],
            "TTFT_p99_ms":  statistics.quantiles(self.ttft_ms, n=100)[98],
            "TPOT_p50_ms":  statistics.median(self.tpot_ms),
            "TPOT_p95_ms":  statistics.quantiles(self.tpot_ms, n=20)[18],
            "throughput_tok_s": sum(self.total_tokens) /
                               (max(self.ttft_ms) / 1000 +
                                sum(self.tpot_ms)/1000),
            "errors": self.errors,
        }

async def benchmark_request(
    session: aiohttp.ClientSession,
    url: str,
    prompt: str,
    max_tokens: int,
    result: BenchmarkResult
):
    payload = {
        "model": "your-model",
        "messages": [{"role": "user", "content": prompt}],
        "max_tokens": max_tokens,
        "stream": True
    }
    start = time.perf_counter()
    first_token_time = None
    token_count = 0

    try:
        async with session.post(f"{url}/v1/chat/completions",
                                json=payload) as resp:
            async for line in resp.content:
                line = line.decode().strip()
                if line.startswith("data: ") and line != "data: [DONE]":
                    chunk = json.loads(line[6:])
                    delta = chunk["choices"][0]["delta"]
                    if delta.get("content"):
                        if first_token_time is None:
                            first_token_time = time.perf_counter()
                            result.ttft_ms.append(
                                (first_token_time - start) * 1000)
                        token_count += 1

        if token_count > 1:
            elapsed = time.perf_counter() - first_token_time
            result.tpot_ms.append(elapsed * 1000 / (token_count - 1))
        result.total_tokens.append(token_count)
    except Exception as e:
        result.errors += 1
        print(f"Error: {e}")

async def run_benchmark(
    url: str = "http://localhost:8000",
    concurrent_users: int = 10,
    num_requests: int = 100,
    prompt: str = "Write a 200-word essay about GPU architecture.",
    max_tokens: int = 300
):
    result = BenchmarkResult()
    connector = aiohttp.TCPConnector(limit=concurrent_users * 2)

    async with aiohttp.ClientSession(connector=connector) as session:
        tasks = [
            benchmark_request(session, url, prompt, max_tokens, result)
            for _ in range(num_requests)
        ]
        # Run with limited concurrency
        semaphore = asyncio.Semaphore(concurrent_users)
        async def bounded_task(t):
            async with semaphore:
                await t
        await asyncio.gather(*[bounded_task(t) for t in tasks])

    print("\n=== Benchmark Results ===")
    for k, v in result.summary().items():
        print(f"  {k:25s}: {v:.2f}" if isinstance(v, float) else
              f"  {k:25s}: {v}")

if __name__ == "__main__":
    asyncio.run(run_benchmark(
        url="http://localhost:8000",
        concurrent_users=20,
        num_requests=200,
    ))
\end{lstlisting}

% ============================================================
\chapter{GPU Profiling Tools}
\label{app:profiling}
% ============================================================

\section*{Profiling Workflow}

\begin{lstlisting}[style=bashstyle, caption={Full GPU profiling workflow}]
# ── Step 1: Quick health check ───────────────────────────────
nvidia-smi
# Look for: GPU util%, memory used, temperature, power draw

# ── Step 2: Find the bottleneck with nsys ────────────────────
nsys profile \
    --trace=cuda,nvtx,osrt,cudnn,cublas \
    --sample=cpu \
    --output=my_profile \
    python my_inference.py

# Open GUI
nsys-ui my_profile.nsys-rep
# KEY QUESTIONS:
# - Are there gaps between kernels? (CPU-bound, batching issue)
# - Is cudaMemcpy dominating? (unnecessary H2D/D2H transfers)
# - Are kernels short? (need to increase batch size)

# ── Step 3: Deep kernel analysis with ncu ────────────────────
# WARNING: ncu is slow. Use on small workloads only.
ncu --set full \
    --kernel-name "cutlass" \    # Target specific kernel
    --launch-skip 10 \           # Skip first 10 launches (warmup)
    --launch-count 3 \           # Profile 3 launches
    -o kernel_analysis \
    python run_matmul.py

ncu-ui kernel_analysis.ncu-rep
# KEY METRICS:
# - Memory Throughput: are we hitting bandwidth ceiling?
# - SM Active Cycles %: is compute busy?
# - L2 Hit Rate: are we reusing data?
# - Warp Stall Reasons: what's blocking warps?

# ── Step 4: MFU (Model FLOP Utilization) ────────────────────
# MFU = actual FLOP/s / peak FLOP/s
# H100 peak FP16: 989 TFLOP/s
# Good MFU for inference: >30% (memory-bound, expected to be low)
# Good MFU for large batches: >50%

# ── Step 5: Memory analysis ──────────────────────────────────
# Check for memory fragmentation
curl http://localhost:8000/metrics | grep cache
# vllm:gpu_cache_usage_perc  -> should be <90%
# vllm:gpu_cache_hit_rate    -> higher is better (prefix caching)
\end{lstlisting}

\section*{Interpreting Common Profiling Issues}

\begin{longtable}{@{}p{4cm}p{4cm}p{5cm}@{}}
  \toprule
  \textbf{Symptom} & \textbf{Likely Cause} & \textbf{Fix} \\
  \midrule
  \endfirsthead
  \toprule
  \textbf{Symptom} & \textbf{Likely Cause} & \textbf{Fix} \\
  \midrule
  \endhead
  GPU util $<$ 30\% & Batch too small & Increase max\_num\_seqs \\
  Memory $>$ 95\% & KV cache overflow & Reduce seq length or add GPUs \\
  TTFT spikes & Large prefill batches & Enable chunked prefill \\
  TPOT high & Memory bandwidth bound & Normal for small batches; use quantization \\
  Frequent OOM & Memory fragmentation & Use PagedAttention (vLLM) \\
  Slow decode after long prefill & KV cache eviction & Increase GPU memory fraction \\
  \bottomrule
  \caption{Common Profiling Issues and Fixes}
\end{longtable}

% ============================================================
\chapter{Research Papers Every Inference Engineer Should Read}
\label{app:papers}
% ============================================================

\section*{Foundational Papers}

\begin{description}
  \item[Attention Is All You Need (Vaswani et al., 2017)]
    The original Transformer paper.  Understand this before anything else.
    \url{https://arxiv.org/abs/1706.03762}

  \item[FlashAttention (Dao et al., 2022)]
    IO-aware exact attention that became standard in all production systems.
    \url{https://arxiv.org/abs/2205.14135}

  \item[FlashAttention-2 (Dao, 2023)]
    Further optimizations: 2$\times$ speedup, better parallelism.
    \url{https://arxiv.org/abs/2307.08691}

  \item[Efficient Memory Management for LLM Serving with PagedAttention
    (Kwon et al., SOSP 2023)]
    Foundation of vLLM.  Must-read for KV cache management.
    \url{https://arxiv.org/abs/2309.06180}
\end{description}

\section*{Quantization}

\begin{description}
  \item[LLM.int8(): 8-bit Matrix Multiplication for Transformers at Scale
    (Dettmers et al., 2022)]
    First practical INT8 quantization for LLMs.
    \url{https://arxiv.org/abs/2208.07339}

  \item[GPTQ: Accurate Post-Training Quantization for Generative Pre-trained
    Transformers (Frantar et al., 2022)]
    State-of-the-art INT4 PTQ.
    \url{https://arxiv.org/abs/2210.17323}

  \item[AWQ: Activation-aware Weight Quantization (Lin et al., 2023)]
    Better accuracy than GPTQ at same bit width.
    \url{https://arxiv.org/abs/2306.00978}

  \item[SmoothQuant (Xiao et al., 2022)]
    Smooth activation outliers for INT8 quantization.
    \url{https://arxiv.org/abs/2211.10438}
\end{description}

\section*{Speculative Decoding}

\begin{description}
  \item[Accelerating Large Language Model Decoding with Speculative Sampling
    (Chen et al., 2023)]
    Original speculative decoding paper.
    \url{https://arxiv.org/abs/2302.01318}

  \item[EAGLE: Speculative Sampling Requires Rethinking Feature Uncertainty
    (Li et al., 2024)]
    3$\times$ speedup with draft head sharing target model features.
    \url{https://arxiv.org/abs/2401.15077}
\end{description}

\section*{Parallelism and Distributed Inference}

\begin{description}
  \item[Megatron-LM (Shoeybi et al., 2019)]
    Efficient tensor and pipeline parallelism for large models.
    \url{https://arxiv.org/abs/1909.08053}

  \item[Alpa: Automating Inter- and Intra-Operator Parallelism (Zheng et al., 2022)]
    Automatic parallelism for large models.
    \url{https://arxiv.org/abs/2201.12023}
\end{description}

\section*{Systems and Scheduling}

\begin{description}
  \item[Orca: A Distributed Serving System for Transformer-Based Generative
    Models (Yu et al., OSDI 2022)]
    First continuous batching paper for LLM serving.
    \url{https://www.usenix.org/conference/osdi22/presentation/yu}

  \item[SGLang: Efficient Execution of Structured Language Model Programs
    (Zheng et al., 2024)]
    RadixAttention for prefix sharing.
    \url{https://arxiv.org/abs/2312.07104}

  \item[DejaVu: Contextual Sparsity for Efficient LLMs at Inference Time
    (Liu et al., 2023)]
    Dynamic sparsity for faster inference.
    \url{https://arxiv.org/abs/2310.17157}
\end{description}

\section*{2025 Frontier Papers}

\begin{description}
  \item[TokenWeave: Reducing Tensor-Parallel Communication Overhead
    (arXiv 2505.11329, 2025)]
    1.2$\times$ latency improvement in tensor-parallel low-latency inference.

  \item[QuantSpec: Self-Speculative Decoding with Hierarchical Quantized KV
    Cache (arXiv 2502.10424, 2025)]
    1.78$\times$ throughput via combined speculative decoding and KV quantization.

  \item[GPU-Accelerated INT8 Quantization for KV Cache Compression
    (arXiv 2601.04719, 2026)]
    4$\times$ memory reduction with CUDA kernels achieving 1694$\times$ CPU speedup.
\end{description}

% ─── Index ───────────────────────────────────────────────────
\printindex

\end{document}
